% PTAPI v2 - documento tecnico.
% Compilar com LuaLaTeX (tres passagens, por causa do sumario):
%     make            (nesta pasta)   ou   latexmk -lualatex documento.tex
\documentclass[10pt,a4paper]{article}
\usepackage{fontspec}
\usepackage{unicode-math}
\setmainfont{Libertinus Serif}
\setsansfont{Libertinus Sans}[BoldItalicFont={Libertinus Sans Bold}, BoldItalicFeatures={FakeSlant=0.18}]
\setmonofont{DejaVu Sans Mono}[Scale=0.80]
\setmathfont{Libertinus Math}
\usepackage[brazilian]{babel}
\usepackage[a4paper,top=24mm,bottom=26mm,left=22mm,right=22mm]{geometry}
\usepackage{microtype}
\usepackage{xcolor}
\usepackage{booktabs,tabularx,array}
\usepackage{enumitem}
\usepackage{titlesec}
\usepackage{listings}
\usepackage{tikz}
\usetikzlibrary{arrows.meta,positioning}
\usepackage[skins,breakable]{tcolorbox}
\usepackage[hidelinks,pdftitle={PTAPI versão 2: alterações, métodos e guia de uso}]{hyperref}

\definecolor{azul}{HTML}{1F3A5F}
\definecolor{ouro}{HTML}{9A8A55}
\definecolor{fundocod}{HTML}{F6F8FA}
\definecolor{fundosai}{HTML}{FBFAF5}
\definecolor{cinza}{HTML}{555555}

\setlength{\parindent}{0pt}
\setlength{\parskip}{0.55em plus 0.1em minus 0.05em}
\linespread{1.06}
\setlist{topsep=0.2em,itemsep=0.15em,parsep=0.1em,leftmargin=1.6em}
\emergencystretch=2em

\titleformat{\section}{\sffamily\bfseries\Large\color{azul}}{\thesection}{0.8em}{}[{\color{azul}\titlerule[0.8pt]}]
\titleformat{\subsection}{\sffamily\bfseries\large\color{azul}}{\thesubsection}{0.8em}{}
\titleformat{\subsubsection}{\sffamily\bfseries\normalsize}{\thesubsubsection}{0.8em}{}
\titlespacing*{\section}{0pt}{2.2em}{0.9em}
\titlespacing*{\subsection}{0pt}{1.5em}{0.4em}
\titlespacing*{\subsubsection}{0pt}{1.1em}{0.2em}

% codigo no meio do texto
\newcommand{\code}[1]{{\ttfamily #1}}

% blocos de codigo (C++) e de saida de programa
\lstdefinestyle{base}{basicstyle=\ttfamily\fontsize{9.3}{10.8}\selectfont, columns=fullflexible, keepspaces=true,
  breaklines=true, breakatwhitespace=false, showstringspaces=false, tabsize=4, extendedchars=true,
  frame=leftline, framerule=1.6pt, framesep=3mm, xleftmargin=3.6mm, aboveskip=0.7em, belowskip=0.5em}
\lstdefinestyle{cpp}{style=base, language=C++, rulecolor=\color{azul}, backgroundcolor=\color{fundocod},
  keywordstyle=\bfseries\color[HTML]{0B5394}, commentstyle=\itshape\color[HTML]{5E7A8A}, stringstyle=\color[HTML]{9C5700},
  morekeywords={override,nullptr,constexpr,noexcept,size_t}}
\lstdefinestyle{saida}{style=base, rulecolor=\color{ouro}, backgroundcolor=\color{fundosai}}

% tabelas
\newcommand{\cab}[1]{\textsf{\textbf{#1}}}
\newenvironment{tabela}[1]{\par\smallskip\small\renewcommand{\arraystretch}{1.18}\tabularx{\linewidth}{@{}#1@{}}}{\endtabularx\par\smallskip}

% caixas
\newtcolorbox{nota}{colback=azul!4, colframe=azul!35, boxrule=0.4pt, arc=1pt, left=3mm, right=3mm, top=1.5mm, bottom=1.5mm,
  fontupper=\small, before skip=0.8em, after skip=0.8em}
\newtcolorbox{fonte}{blanker, breakable, borderline west={1.4pt}{0pt}{ouro}, left=3.2mm, fontupper=\small\color{black!80},
  before skip=0.5em, after skip=0.8em}
\newenvironment{referencias}{\small\newcommand{\refitem}{\par\hangindent=7mm\hangafter=1 }}{\par}

% Figura 1: dependencias entre os nos
\newcommand{\FiguraGrafo}{%
\begin{center}
\begin{tikzpicture}[x=1mm,y=1mm, font=\sffamily\footnotesize, >={Stealth[length=2mm]},
  mc/.style={draw=azul, fill=azul!10, rounded corners=1pt, minimum width=24mm, minimum height=8mm},
  sw/.style={draw=ouro!80!black, fill=ouro!22, rounded corners=1pt, minimum width=24mm, minimum height=9mm, align=center}]
  \foreach \i in {0,1,2,3}{
    \node[mc] (a\i) at (40*\i, 44) {MC\textsubscript{\i} ($T_{\i}$)};
    \node[mc] (b\i) at (40*\i, 0)  {MC\textsubscript{\i} ($T_{\i}$)};
  }
  \foreach \i in {0,1,2}{
    \node[sw] (s\i) at (40*\i+20, 22) {SW\textsubscript{\i}\\[-1pt]{\scriptsize\color{cinza}(+ TU\textsubscript{\i})}};
  }
  \foreach \i/\j in {0/1,1/2,2/3}{
    \draw[->] (a\i.south) -- (s\i);
    \draw[->] (a\j.south) -- (s\i);
    \draw[->] (s\i) -- (b\i.north);
    \draw[->] (s\i) -- (b\j.north);
  }
  \draw[->] (s0) -- (s1);
  \draw[->] (s1) -- (s2);
  \node[anchor=east, text=cinza] at (-14, 44) {rodada $r$};
  \node[anchor=east, text=cinza] at (-14, 0) {rodada $r+1$};
\end{tikzpicture}
\end{center}
{\small\sffamily\color{black!75}\textbf{Figura 1.} Dependências entre os nós para quatro temperaturas. Cada MC executa MCL passos de
Metropolis; cada SW propõe a troca das réplicas do seu par e, a cada $\upsilon$ propostas, executa a atualização de temperatura (TU).
A cadeia SW\textsubscript{0} $\to$ SW\textsubscript{1} $\to$ SW\textsubscript{2} fixa a ordem das trocas dentro da rodada.\par}
\medskip
}

\begin{document}

\begin{titlepage}
\sffamily
\vspace*{70mm}
{\fontsize{32}{36}\selectfont\bfseries\color{azul} PTAPI — versão 2\par}
\vspace{6mm}
{\Large Alterações realizadas, origem dos novos métodos\\[2pt] e guia detalhado de funcionamento e uso\par}
\vspace{18mm}
{\color{cinza}
Parallel Tempering API em C++ (modelo \textit{dataflow} para CPU)\\[3pt]
Documento técnico da versão candidata 2.0 — novembro de 2026\\[3pt]
Código descrito: este repositório (núcleo em \code{include/})\par}
\end{titlepage}

\tableofcontents
\clearpage

\section{O que é este documento}\label{s1}

Este documento registra tudo o que mudou da PTAPI original (aqui chamada de v1; o núcleo dela está na pasta \code{v1/\allowbreak{}}) para a versão 2, explica como a v2 funciona por dentro e mostra, com programas completos, como usá-la. Para cada método novo é indicada a publicação de onde ele veio; quando uma regra é adaptação ou procedimento próprio desta versão, isso é dito explicitamente, para que não se atribua a um artigo o que ele não contém.

A v2 foi construída sob quatro princípios:

\begin{itemize}
\item \textbf{A API continua genérica.} O núcleo não conhece o problema; um problema novo só precisa de uma estrutura de solução e de uma classe derivada de \code{Problem<S>}.
\item \textbf{Compatibilidade de código.} Programas escritos para a v1 compilam sem alteração: o construtor de \code{PT}, \code{start()}, \code{getBestSol()}, \code{getBestPTLSolIndex()}, \code{getPTLEnd()}, \code{setTimeEnd()} e \code{setPTLStopEnd()} não mudaram.
\item \textbf{Toda alteração é medida.} Tempo (comparação com a v1 no mesmo programa) e qualidade (média de 20 sementes) foram medidos a cada mudança; os números estão na \hyperref[s9]{Seção 9}.
\item \textbf{Testes unitários permanentes.} A suíte de \code{tests/\allowbreak{}} tem 44 testes e passa em três modos: normal, ThreadSanitizer e AddressSanitizer + UBSan. \end{itemize}

O código da v1 não foi modificado. As medições deste documento foram feitas em uma máquina de 24 núcleos lógicos com g++ 15.2 e \code{-O3}.

\begin{nota}
\textbf{Convenção.} $T$ é a temperatura e $\beta = 1/T$. As temperaturas são numeradas de $T_{0}$ (a primeira) a $T_{M-1}$ (a última); nas distribuições 1, 3 e 4 a primeira é a mais fria. $E$ é o valor da função objetivo (\code{evalSol}); a API minimiza. MCL é o comprimento de cada cadeia de Markov e PTL o número de rodadas de troca.
\end{nota}

\section{Como a API v2 funciona}\label{s2}

\subsection{Os três passos}

O uso da API tem os mesmos três passos da v1:

\begin{enumerate}
\item \textbf{Solução}: uma \code{struct} que herda de \code{solution}. A base traz o campo \code{double evalSol}, onde fica o custo; o restante é a codificação do problema.
\item \textbf{Problema}: uma classe que herda de \code{Problem<S>} e implementa \code{construction()} (solução inicial avaliada), \code{neighbor()} (um vizinho) e \code{evaluate()} (o custo).
\item \textbf{PT}: um objeto \code{PT<S>} criado com os parâmetros do algoritmo, e a chamada \code{start(\allowbreak{}threads, \&problema)}, que devolve a melhor solução. \end{enumerate}

O construtor tem oito argumentos, os mesmos da v1:

\begin{lstlisting}[style=cpp]
PT<S> pt(tempMin, tempMax, replicas, MCL, PTL,
         distribuicao, tipoDeAtualizacao, intervaloDeAtualizacao);
\end{lstlisting}

\begin{tabela}{l>{\raggedright\arraybackslash}X}
\toprule
\cab{Argumento} & \cab{Significado} \\
\midrule
\code{tempMin}, \code{tempMax} & Primeira e última temperatura da escada. \\
\code{replicas} & Número de temperaturas (uma réplica em cada). \\
\code{MCL} & Movimentos que cada réplica faz entre duas propostas de troca. \\
\code{PTL} & Número de rodadas (cada rodada: MCL movimentos por réplica e as propostas de troca). \\
\code{distribuicao} & Distribuição inicial das temperaturas: 1 linear, 2 linear inversa, 3 exponencial, 4 progressão geométrica. \\
\code{tipoDeAtualizacao} & Regra de atualização dinâmica das temperaturas, de 0 a 6 (\hyperref[s7]{Seção 7}). \\
\code{intervaloDeAtualizacao} & Número de propostas de troca entre duas atualizações das temperaturas. \\
\bottomrule
\end{tabela}

\subsection{O grafo \textit{dataflow}}

A execução é descrita por um grafo de tarefas. Há três tipos de nó:

\begin{itemize}
\item \textbf{MC} (\code{NodeMCMC}): uma cadeia de Markov homogênea de MCL passos em uma temperatura. Há um nó MC por temperatura; ele guarda a réplica que está naquela temperatura.
\item \textbf{SW} (\code{NodeSwap}): a proposta de troca entre duas temperaturas adjacentes. Há um nó SW por par.
\item \textbf{TU} (\code{NodeTempUp}): a atualização das temperaturas de um par. Não é escalonado separadamente: é executado dentro do nó SW a que pertence, a cada \code{intervaloDeAtualizacao} propostas. \end{itemize}

As arestas dizem quem espera por quem. O nó SW\textsubscript{$i$} só pode executar depois que MC\textsubscript{$i$} e MC\textsubscript{$i+1$} terminaram a cadeia da rodada, e depois que SW\textsubscript{$i-1$} executou. Os nós MC\textsubscript{$i$} da rodada seguinte ficam prontos quando as trocas que os envolvem (SW\textsubscript{$i-1$} e SW\textsubscript{$i$}) terminaram. Os nós são reutilizados a cada rodada; a figura desdobra duas rodadas para mostrar as dependências.

\FiguraGrafo

Essa estrutura é a mesma da tese; o que mudou na v2 é a forma de executá-la (\hyperref[s4]{Seção 4}). Com uma única réplica o grafo é um nó MC ligado a si mesmo: Metropolis simples.

\subsection{Uma rodada}

\textbf{Nó MC.} Para a solução corrente $s$ na temperatura $T$, repete MCL vezes: gera um vizinho avaliado $s'$, calcula $\Delta = E(s') - E(s)$ e aceita o vizinho se

\[
\Delta \le 0 \quad\text{ou}\quad u \le \exp(-\Delta/T), \qquad \text{com } u \text{ uniforme em } [0,1).
\]

Quando um movimento melhora a solução da réplica, ela é comparada com a melhor solução global, que é atualizada na hora (na v1 a melhor era coletada dos nós só no final).

\textbf{Nó SW.} Para as réplicas das temperaturas $T_{i}$ e $T_{i+1}$, com custos $E_{i}$ e $E_{i+1}$, a troca é aceita com probabilidade

\[
p = \min\bigl\{1,\ \exp\bigl[(\beta_{i+1}-\beta_i)(E_{i+1}-E_i)\bigr]\bigr\}.
\]

Trocar significa trocar as soluções entre os dois nós MC (um \code{std::swap}, sem cópia); as temperaturas ficam nos nós. Junto com a solução viaja um rótulo usado para medir idas e voltas: a réplica recebe o rótulo «subindo» quando passa pela primeira temperatura e «descendo» quando passa pela última.

\textbf{Fim.} A execução termina quando todas as réplicas completaram PTL rodadas ou quando um dos critérios opcionais é atingido (\hyperref[s6]{Seção 6.3}).

\subsection{O \textit{pool} de threads}

Os nós prontos ficam em uma única fila. Cada thread do \textit{pool} retira um nó, executa-o e avisa os sucessores; cada nó tem um contador atômico de dependências pendentes e entra na fila quando o contador chega a zero. Três decisões reduzem o custo do escalonamento:

\begin{itemize}
\item \textbf{Continuação na própria thread.} Ao terminar um nó, a thread executa ela mesma o próximo sucessor que ficou pronto, sem passar pela fila. O sucessor preferido de um SW é o SW seguinte, de modo que uma thread percorre a cadeia de trocas (que é barata) e entrega os nós MC às outras.
\item \textbf{Espera em duas fases.} Uma thread sem trabalho gira por um número limitado de iterações (\code{PTAPI\_\allowbreak{}SPIN}, padrão 5000) e só então dorme. Na v1 as threads ociosas ficavam em espera ocupada, a 100\% de CPU, inclusive a thread que chamou \code{start()}.
\item \textbf{Número de threads.} São criadas min(\code{thN}, réplicas) threads, porque nunca há mais nós MC prontos do que réplicas; \code{thN $\le$ 0} significa todos os núcleos. \end{itemize}

O valor 5000 de \code{PTAPI\_\allowbreak{}SPIN} foi escolhido medindo de 0 a 200000 em máquina ociosa e sob carga; pode ser alterado na compilação com \code{-DPTAPI\_\allowbreak{}SPIN=n}.

\subsection{Números aleatórios e reprodutibilidade}

Cada nó (MC e SW) tem o seu próprio gerador \code{std::mt19937\_\allowbreak{}64}, derivado de uma única semente e do índice do nó. O código do problema obtém o gerador da réplica que está executando com \code{ptapi::rng()}, \code{ptapi::randInt(\allowbreak{}a, b)} e \code{ptapi::rand01()}. Na v1 todos os nós eram semeados com a mesma semente padrão, e os exemplos criavam um \code{std::random\_\allowbreak{}device} e um gerador novo a cada chamada de \code{neighbor()}, o que é lento e impede a repetição de uma execução.

Como o grafo fixa a ordem em que os dados são lidos e escritos, uma execução que só usa os geradores da API e termina por PTL dá o mesmo resultado para qualquer número de threads. Isso vale para os tipos de atualização 0, 1, 2 e 5; os tipos 3, 4 e 6 leem estatísticas de outros pares enquanto eles executam e por isso não são garantidamente reprodutíveis (\hyperref[s7]{Seção 7}). Paradas por tempo também não são reprodutíveis, por natureza.

\subsection{Onde termina a API e começa o problema}

A classe do problema declara \code{evaluate} em qualquer uma de três formas; a API descobre qual em tempo de compilação, a partir do tipo do ponteiro passado a \code{start()} (por isso o ponteiro deve ser da própria classe, \code{SSP*}, e não \code{Problem<solSSP>*}):

\begin{lstlisting}[style=cpp]
S      evaluate(S sol);          // devolve a solução avaliada (pode também alterá-la)
double evaluate(S sol);          // forma dos exemplos publicados com a v1
double evaluate(const S& sol);   // sem cópia da solução
\end{lstlisting}

O passo da cadeia chama uma única função virtual, \code{neighborEval(\allowbreak{}const S\& atual, S\& candidata)}, cuja implementação padrão é \code{candidata = evaluate(\allowbreak{}neighbor(\allowbreak{}atual))}. Sobrescrevê-la é opcional e é o que mais acelera um problema (\hyperref[s6]{Seção 6.1}).

\clearpage
\section{Guia de uso com um exemplo completo}\label{s3}

O programa abaixo resolve a partição de números: dividir 60 inteiros em dois grupos cujas somas sejam tão próximas quanto possível. Ele é o arquivo \code{docs/\allowbreak{}report/\allowbreak{}src/\allowbreak{}guia.cpp}, compilado e executado para gerar a saída mostrada. Foi escolhido um problema diferente dos exemplos da API para mostrar que só a classe do problema muda.

\subsubsection*{Passos 1 e 2: solução e problema}

\lstinputlisting[style=cpp,firstline=4,lastline=59]{guia.cpp}

Pontos a observar:

\begin{itemize}
\item \code{construction()} devolve a solução já avaliada (\code{evalSol} preenchido).
\item Os números aleatórios vêm de \code{ptapi::randInt}; é isso que torna a execução reprodutível com \code{setSeed()}.
\item \code{neighbor()} e \code{evaluate()} bastam para a API funcionar. \code{neighborEval()} é a versão rápida do mesmo passo: reaproveita a memória de \code{cand} e calcula o custo a partir do custo atual, em tempo constante. A solução guarda o dado auxiliar \code{diff} exatamente para permitir essa avaliação incremental.
\item Na entrada de \code{neighborEval}, \code{cand} é uma solução válida de conteúdo qualquer (em geral a solução rejeitada no passo anterior); por isso todos os seus campos devem ser escritos. \end{itemize}

\subsubsection*{Passo 3: o PT}

\lstinputlisting[style=cpp,firstline=61,lastline=76]{guia.cpp}

Compilação e execução (C++17 ou mais novo; a API é só de cabeçalhos):

\begin{lstlisting}[style=saida]
$ g++ guia.cpp -std=c++17 -O3 -pthread -o guia
$ ./guia
melhor diferenca: 1 (rodada 1411 de 2000, 97.8 ms de 138.2 ms)
12800000 avaliacoes em 16 threads, semente 42
idas e voltas: 12   barreira: 5.96
escada final : 10 22 46 100 215 464 1000 2154 4642 10000 21544 46416 100000 215443 464159 1000000
taxa de troca: 0.40 0.47 0.51 0.59 0.65 0.71 0.62 0.66 0.62 0.63 0.65 0.61 0.62 0.64 0.65
\end{lstlisting}

\subsubsection*{Lendo a saída}

\begin{itemize}
\item \textbf{Melhor solução}: diferença 1 entre os grupos, encontrada na rodada 1411. As 2000 rodadas de 16 réplicas com MCL 400 somam 12,8 milhões de avaliações em 138 ms.
\item \textbf{Taxa de troca}: há um valor por par de temperaturas adjacentes. Taxas entre 40\% e 71\% indicam uma escada mais densa que o necessário: 16 réplicas são mais do que o intervalo de 10 a $10^{6}$ pede neste problema.
\item \textbf{Barreira}: a soma das taxas de rejeição das trocas, 5,96 (\hyperref[s6]{Seção 6.7}).
\item \textbf{Idas e voltas}: 12 réplicas foram da primeira temperatura à última e voltaram. \end{itemize}

\subsubsection*{Deixando a API ajustar a escada}

Trocando apenas o sétimo argumento do construtor de 0 para 5 (aproximação estocástica, \hyperref[s7]{Seção 7}), a mesma execução termina com esta escada:

\begin{lstlisting}[style=saida]
melhor diferenca: 1 (rodada 71 de 2000, 5.6 ms de 163.1 ms)
idas e voltas: 10   barreira: 4.97
escada final : 10 19 118 1105 7668 53993 414053 1000000 1000000 1000000 1000000 1000000 ...
taxa de troca: 0.33 0.30 0.28 0.27 0.26 0.24 0.55 0.89 0.95 0.99 1.00 1.00 1.00 0.99 1.00
\end{lstlisting}

A regra espaçou as temperaturas até as taxas ficarem perto de 23\% e, com isso, a escada chegou à temperatura máxima na oitava réplica; as outras oito ficaram empilhadas no limite. Essa é a leitura prática do resultado: para este problema e este intervalo, cerca de oito réplicas bastam. Os dois resultados (tipo 0 e tipo 5) repetem-se exatamente em outra execução com a mesma semente.

\section{Alterações internas e desempenho}\label{s4}

O algoritmo e o grafo são os da v1. O que mudou foi a implementação do escalonamento e do passo da cadeia:

\begin{tabela}{>{\raggedright\arraybackslash}p{0.22\linewidth}>{\raggedright\arraybackslash}p{0.36\linewidth}>{\raggedright\arraybackslash}X}
\toprule
\cab{Tema} & \cab{v1} & \cab{v2} \\
\midrule
Threads ociosas & Espera ocupada (100\% de CPU), inclusive a thread que chamou \code{start()}. & Giram por alguns microssegundos e depois dormem. \\
Número de threads & Sempre \code{thN}. & min(\code{thN}, réplicas); \code{thN $\le$ 0} usa todos os núcleos. \\
Dependências dos nós & Mutex e vetor de marcas percorrido a cada notificação. & Um contador atômico por nó. \\
Filas & Uma fila com mutex por thread, percorridas em laço. & Uma fila; a thread executa ela mesma o próximo nó pronto. \\
Cópias por passo MC & 3 a 4 cópias profundas da solução. & 1 cópia (nenhuma com \code{neighborEval}). \\
Troca de réplicas & 3 a 4 cópias profundas. & \code{std::swap}. \\
Melhor solução & Coletada dos nós ao final. & Global, atualizada quando é encontrada. \\
Números aleatórios & Todos os nós com a mesma semente padrão. & Um gerador por nó, derivado de uma semente. \\
Memória & Nós e \textit{pool} nunca liberados. & Liberados ao fim de \code{start()}. \\
Exceções do problema & Encerravam o programa. & Param a execução e são relançadas por \code{start()}. \\
\bottomrule
\end{tabela}

\subsection*{Tempo medido}

O programa \code{bench/\allowbreak{}bench.cpp} usa só a interface comum às duas versões e é compilado uma vez com cada uma. A tabela traz o tempo de parede em máquina ociosa (menor é melhor). «TSP» é um problema sintético de avaliação barata, que expõe o custo da própria API; «SD» é o problema real da pasta \code{v1/\allowbreak{}SD/\allowbreak{}}, de avaliação cara.

\begin{tabela}{>{\raggedright\arraybackslash}Xrrr}
\toprule
\cab{Problema [réplicas, MCL, PTL, threads]} & \cab{v1 (s)} & \cab{v2 (s)} & \cab{v2 / v1} \\
\midrule
TSP 100 [8, 500, 2000, 8] & 0,702 & 0,604 & 0,86 \\
TSP 100 [8, 500, 2000, 24] & 0,814 & 0,596 & 0,73 \\
TSP 100 [30, 500, 500, 24] & 0,323 & 0,276 & 0,85 \\
TSP 100 [16, 20, 20000, 16] & 0,507 & 0,471 & 0,93 \\
SD [8, 300, 60, 8] & 3,76 & 3,79 & 1,01 \\
SD [8, 300, 60, 24] & 4,60 & 3,99 & 0,87 \\
SD [30, 300, 16, 24] & 1,76 & 1,64 & 0,93 \\
\bottomrule
\end{tabela}

No caso SD [8, 300, 60, 8] o tempo de parede é igual, mas o tempo de CPU caiu de 33,8 s para 25,9 s: a diferença é a espera ocupada que deixou de existir. As quatro linhas de TSP foram medidas com a versão atual, depois de todas as alterações descritas neste documento (médias de 6 a 8 execuções intercaladas); as de SD foram medidas antes da inclusão de checkpoint, movimentos, esquema de troca e calibração e não foram repetidas. Nenhuma das funções acrescentadas depois da primeira medição alterou o laço das cadeias, e a regra do projeto é que a v2 não pode ser mais lenta que a v1: o benchmark é repetido a cada bloco de alterações.

Esses números usam a interface da v1 (\code{neighbor} + \code{evaluate}). O ganho maior está em \code{neighborEval}, que depende do problema: no TSP de 200 cidades a mesma busca, com a mesma semente e o mesmo resultado, passou de 2,6–2,9 para 13,2 milhões de avaliações por segundo (\hyperref[s6]{Seção 6.1}).

\section{Comportamentos diferentes da v1}\label{s5}

O código da v1 compila sem alteração, mas alguns resultados numéricos mudam. As diferenças são as seguintes; \code{setV1Compatibility(\allowbreak{}true)} restaura as três primeiras, para reproduzir resultados já publicados.

\begin{itemize}
\item \textbf{Critério de troca com objetivo real.} A v1 convertia a energia das réplicas para \code{int} antes do teste de troca (e guardava o melhor valor como \code{int}). A v2 usa \code{double}. Nada muda para objetivos inteiros.
\item \textbf{Taxa de troca usada pelo tipo 1 (23\%).} A v1 dividia as trocas aceitas desde a última atualização pelo número total de rodadas; a v2 divide pelo tamanho da janela, como descrito na tese.
\item \textbf{Regras dos tipos 1 e 2.} Na v2 elas mantêm a primeira e a última temperatura e a escada ordenada (\hyperref[s7]{Seção 7}). Na v1 o tipo 1 só aproximava temperaturas, de modo que a escada descia e a temperatura máxima se perdia; o tipo 2 aceitava qualquer valor positivo.
\item \textbf{Tipo 3.} A regra é a da tese, sem alteração; apenas os valores compartilhados entre pares passaram a ser atômicos (havia uma condição de corrida apontada pelo ThreadSanitizer).
\item \textbf{Rótulos de ida e volta.} São guardados pela API. Os campos \code{Nup}/\code{Ndown} continuam em \code{solution} só para que o código antigo compile; são ignorados.
\item \textbf{Limite de tempo.} É verificado pela thread que chamou \code{start()}; os nós MC param no passo seguinte, e não depois de mais uma ou duas rodadas.
\item \textbf{Parâmetros inválidos} (distribuição, tipo de atualização, zero réplicas) lançam \code{std::invalid\_\allowbreak{}argument}. Uma réplica é aceita.
\item \code{start()} pode ser chamado mais de uma vez no mesmo objeto, e os cabeçalhos podem ser incluídos em mais de um \code{.cpp}.
\item \code{using namespace std;} continua exportado por \code{PT.h}, por compatibilidade; \code{-DPTAPI\_\allowbreak{}NO\_\allowbreak{}USING\_\allowbreak{}NAMESPACE\_\allowbreak{}STD} desliga. \end{itemize}

O efeito da regra antiga do tipo 1 aparece no exemplo \code{exampleV1Compatibility} (escada geométrica de 1 a 60, 12 réplicas, 2000 rodadas):

\begin{lstlisting}[style=saida]
tipo 1, regra da v2:  escada final  1.00 1.17 1.20 1.88 2.16 3.44 6.28 9.19 13.95 21.61 41.16 60.00
tipo 1, regra da v1:  escada final  1.00 1.02 1.00 0.97 0.89 0.83 0.81 0.79  0.75  0.62  0.53  0.52
\end{lstlisting}

\section{Novas funções}\label{s6}

Todas são opcionais. As de configuração são chamadas antes de \code{start()}; as de resultado descrevem o último \code{start()}. Para cada grupo há um programa curto em \code{example/\allowbreak{}functions/\allowbreak{}}, sobre o mesmo problema (TSP euclidiano aleatório com movimento 2-opt, em \code{tsp.h}); as saídas reproduzidas abaixo são desses programas.

\subsection{Avaliação e passo da cadeia}

\begin{lstlisting}[style=cpp]
virtual void neighborEval(const S& atual, S& candidata);   // em Problem<S>
\end{lstlisting}

Escreve em \code{candidata} um vizinho avaliado de \code{atual}. A implementação padrão chama \code{neighbor()} e \code{evaluate()}. Sobrescrevê-la permite (a) reaproveitar a memória já alocada em \code{candidata}, eliminando as cópias, e (b) calcular o custo de forma incremental. No TSP com 2-opt, por exemplo, o custo novo sai de quatro distâncias:

\begin{lstlisting}[style=cpp]
void neighborEval(const solTSP& cur, solTSP& cand){
    int a, b;
    do{ a = ptapi::randInt(0, n-1); b = ptapi::randInt(0, n-1); }while(a == b);
    if(a > b) std::swap(a, b);

    cand.sol = cur.sol;                                   // reaproveita a memória de 'cand'
    std::reverse(cand.sol.begin()+a, cand.sol.begin()+b+1);

    if(a == 0 && b == n-1){ cand.evalSol = cur.evalSol; return; }
    int pa = cur.sol[(a+n-1)%n], ca = cur.sol[a], cb = cur.sol[b], nb = cur.sol[(b+1)%n];
    cand.evalSol = cur.evalSol - dist(pa,ca) - dist(cb,nb) + dist(pa,cb) + dist(ca,nb);
}
\end{lstlisting}

Saída de \code{exampleEvaluate} (200 cidades, mesma semente nas quatro execuções):

\begin{lstlisting}[style=saida]
S evaluate(S)                    best  11596     705.6 ms     2.6 million evaluations/s
double evaluate(S)               best  11596     628.5 ms     2.9 million evaluations/s
double evaluate(const S&)        best  11596     696.6 ms     2.6 million evaluations/s
neighborEval() with delta        best  11596     136.1 ms    13.2 million evaluations/s
the four forms reached the same solution
\end{lstlisting}

Se a classe sobrescreve \code{neighborEval} e não são dadas soluções iniciais, \code{evaluate} nem precisa existir.

\subsection{Semente}

\begin{tabela}{>{\raggedright\arraybackslash}p{0.36\linewidth}>{\raggedright\arraybackslash}X}
\toprule
\cab{Função} & \cab{Efeito} \\
\midrule
\code{setSeed(\allowbreak{}unsigned long long)} & Define a semente da execução. \\
\code{getSeed()} & Semente usada no último \code{start()}. Sem \code{setSeed} a API sorteia uma; lê-la permite repetir depois uma execução interessante. \\
\code{ptapi::rng()}, \code{ptapi::randInt(\allowbreak{}a, b)}, \code{ptapi::rand01()} & Gerador da réplica em execução, inteiro uniforme em [a, b] e real uniforme em [0, 1), para uso dentro da classe do problema. \\
\bottomrule
\end{tabela}

\begin{lstlisting}[style=saida]
seed 42   1 thread: 7518   4 threads: 7518   all cores (12): 7518
seed 43   7535
drawn seed 8806199900061321338   7552   repeated: 7552
\end{lstlisting}

\subsection{Critérios de parada}

O PTL do construtor é sempre um limite. Os demais são opcionais e podem ser combinados; o primeiro que for atingido encerra a execução.

\begin{tabela}{>{\raggedright\arraybackslash}p{0.36\linewidth}>{\raggedright\arraybackslash}X}
\toprule
\cab{Função} & \cab{Efeito} \\
\midrule
\code{setTimeLimit(\allowbreak{}double segundos)} & Limite de tempo com fração de segundo (\code{setTimeEnd(\allowbreak{}int)} da v1 continua aceito). \\
\code{setTargetValue(\allowbreak{}double)} & Para quando a melhor solução for menor ou igual ao alvo. \\
\code{setPTLStopEnd(\allowbreak{}float fração)} & Da v1: para depois dessa fração do PTL sem nova melhor solução. \\
\code{setOnImprove(\allowbreak{}f)} & \code{f(\allowbreak{}const S\& sol, int rodada, double ms)} é chamada a cada nova melhor solução, de uma thread de trabalho, uma chamada por vez. \\
\code{requestStop()} & Encerra um \code{start()} em andamento; pode ser chamada de outra thread ou de dentro de \code{setOnImprove}. \\
\code{setStagnationRounds(int)} & Para depois desse número de rodadas sem nova melhor solução (número absoluto, e não fração do PTL). \\
\code{setStagnationTime(double segundos)} & Para depois desse tempo sem nova melhor solução. \\
\code{setLowerBound(limite, folga)} & Para quando a melhor solução estiver a \code{folga} (relativa) de um limite inferior conhecido: melhor $\le$ limite $+$ folga $\times$ |limite|. \\
\code{setStopCondition(f)} & Qualquer outra regra: \code{f(const PTStatus\&)} é chamada ao fim de cada rodada e encerra a execução quando devolve \code{true}. \code{PTStatus} traz \code{round}, \code{elapsedMs}, \code{bestValue}, \code{bestRound} e \code{bestMs}. \\
\code{getStopReason()}, \code{getStopReasonName()} & Por que a execução terminou: PTL, tempo, alvo, estagnação, pedido do usuário, condição ou erro. \\
\bottomrule
\end{tabela}

\begin{lstlisting}[style=cpp]
PT<solTSP> algo(1.0f, 60.0f, 12, 200, 2000000000, 4, 0, 50);   // PTL "sem limite"
algo.setTimeLimit(10.0);                                       // rede de segurança
algo.setTargetValue(9000.0);                                   // o que realmente para
algo.setOnImprove([&](const solTSP& sol, int rodada, double ms){
    printf("novo melhor %.0f na rodada %d (%.1f ms)\n", sol.evalSol, rodada, ms);
});
solTSP s = algo.start(0, &prob);
\end{lstlisting}

\begin{lstlisting}[style=saida]
time limit of 0.25 s         best   7468   stopped at round    1763     250.9 ms
target value of 9000         best   8999   stopped at round      44       6.1 ms
10% of PTL without improving best   7455   stopped at round    3246     509.9 ms
stopped by another thread    best   7459   stopped at round    1445     151.4 ms
\end{lstlisting}

Um critério escrito pelo usuário e, em seguida, a saída de \code{exampleStop} para os gatilhos novos; a última coluna é o motivo do término:

\begin{lstlisting}[style=cpp]
// para quando a busca passou, sem melhorar, três vezes o tempo da melhor atual (no máximo 0,5 s)
algo.setStopCondition([](const PTStatus& st){
    return st.elapsedMs > 50.0 && (st.elapsedMs - st.bestMs) > std::min(3.0*st.bestMs, 500.0);
});
\end{lstlisting}

\begin{lstlisting}[style=saida]
800 rounds without improving   best   7533   round    1969      58.8 ms   stop: stagnation
0.15 s without improving       best   7386   round   13109     370.2 ms   stop: stagnation
within 5% of the bound 7300    best   7645   round     214       6.4 ms   stop: target
user-defined condition         best   7374   round   40358    1166.7 ms   stop: condition
\end{lstlisting}

\subsection{Escada de temperaturas e soluções iniciais}

\begin{tabela}{>{\raggedright\arraybackslash}p{0.36\linewidth}>{\raggedright\arraybackslash}X}
\toprule
\cab{Função} & \cab{Efeito} \\
\midrule
\code{setTemperatures(\allowbreak{}vector<double>)} & Escada definida pelo usuário: quaisquer valores e qualquer número de réplicas (substitui as temperaturas e o número de réplicas do construtor). \\
\code{getInitialTemperatures()} & Escada inicial. \\
\code{getTemperatures()} & Escada final, depois das atualizações dinâmicas. \\
\code{setInitialSolutions(\allowbreak{}vector<S>)} & Partida a quente: as réplicas começam de soluções dadas, em vez de \code{construction()}. A API avalia as soluções; se houver menos soluções que réplicas, elas são atribuídas ciclicamente. \\
\code{getReplicas()} & Solução final de cada temperatura; serve de entrada para uma segunda etapa. \\
\bottomrule
\end{tabela}

As quatro distribuições do construtor, de 1 a 60 com 8 réplicas (a distribuição 2 vai da máxima para a mínima, como na v1; as outras são crescentes):

\begin{lstlisting}[style=saida]
1 linear            1.00    9.43   17.86   26.29   34.71   43.14   51.57   60.00
2 inverse linear   60.00    6.36    3.36    2.28    1.73    1.39    1.16    1.00
3 exponential       1.00    1.79    3.22    5.78   10.38   18.63   33.43   60.00
4 geometric         1.00    1.79    3.22    5.78   10.38   18.63   33.43   60.00
\end{lstlisting}

Partida a quente no TSP de 200 cidades, em uma execução curta (\code{exampleInitialSolutions}):

\begin{lstlisting}[style=saida]
nearest neighbour tour (city 0)   13908
random start                      13651   first best at 106246
start from nearest neighbour      11443   first best at 13908
second stage from its replicas    11170
\end{lstlisting}

\textbf{Heurística construtiva dentro do problema.} A função \code{constructionReplica(indice, total)} pode ser sobrescrita na classe do problema; ela constrói a solução de uma réplica sabendo a sua posição na escada (0 é a primeira temperatura). A implementação padrão chama \code{construction()}. Assim a heurística faz parte do problema e nada precisa ser preparado no programa principal:

\begin{lstlisting}[style=cpp]
// metade fria da escada: vizinho mais próximo, uma cidade inicial por réplica;
// metade quente: soluções aleatórias, que mantêm a diversidade
solTSP constructionReplica(int index, int count) override {
    if(index >= count/2) return construction();
    solTSP s = nearestNeighbour(index % n);
    s.evalSol = evaluate(s);                 // deve devolver a solução avaliada
    return s;
}
\end{lstlisting}

Com \code{setParallelConstruction(true)} as réplicas são construídas por várias threads ao mesmo tempo, o que compensa quando a heurística é cara: no teste da suíte, 8 réplicas com construção de 30 ms cada levaram 241 ms em sequência e 30 ms em paralelo. As soluções construídas não dependem das threads, porque cada réplica tem o seu gerador aleatório; o código de construção precisa apenas ser seguro para chamadas simultâneas, como \code{neighbor()} já é. O padrão continua sendo a construção em sequência, como na v1.

\subsection{Vários movimentos com probabilidades}

Um problema com mais de um tipo de movimento registra um peso para cada um e pergunta à API qual aplicar. Depois da execução, a API informa quantas vezes cada movimento foi proposto e aceito.

\begin{lstlisting}[style=cpp]
class TSPMoves: public Problem<solTSP>{
  public:
    TSPMoves(int n_): base(n_), n(n_){
        setMoveProbabilities({0.6, 0.1, 0.3});       // pesos; a API normaliza
    }
    solTSP neighbor(solTSP s){
        int a, b;
        do{ a = ptapi::randInt(0, n-1); b = ptapi::randInt(0, n-1); }while(a == b);
        switch(chooseMove()){                        // sorteia o movimento
            case 0:  /* 2-opt     */ if(a > b) std::swap(a, b);
                                     std::reverse(s.sol.begin()+a, s.sol.begin()+b+1); break;
            case 1:  /* troca     */ std::swap(s.sol[a], s.sol[b]); break;
            default: /* inserção  */ { int c = s.sol[a]; s.sol.erase(s.sol.begin()+a);
                                       s.sol.insert(s.sol.begin()+b, c); }
        }
        return s;
    }
    // ...
};

for(const PTMoveStat& m : algo.getMoveStatistics())   // {proposed, accepted}, um por movimento
    printf("%lld propostos, %lld aceitos\n", m.proposed, m.accepted);
\end{lstlisting}

\begin{lstlisting}[style=saida]
60% 2-opt, 10% swap, 30% insertion: best 7510
   2-opt      probability 0.60   proposed   1440556   accepted    17578 ( 1.2%)
   swap       probability 0.10   proposed    239866   accepted     1937 ( 0.8%)
   insertion  probability 0.30   proposed    719578   accepted     7418 ( 1.0%)
\end{lstlisting}

As contagens somam todas as temperaturas. Se a classe não registra movimentos, o laço das cadeias é o mesmo de antes e nada é contado. As probabilidades são fixas durante a execução; a API não as adapta.

\subsection{Checkpoint e retomada}

O arquivo de checkpoint guarda as temperaturas, a solução de cada réplica, a melhor solução, e as rodadas e o tempo já usados. A API escreve os valores da função objetivo; o usuário fornece duas funções, uma que escreve e outra que lê o restante de uma solução.

\begin{lstlisting}[style=cpp]
auto escreve = [](std::ostream& out, const solTSP& s){ for(int c : s.sol) out << c << ' '; };
auto le      = [n](std::istream& in, solTSP& s){ s.sol.resize(n); for(int& c : s.sol) in >> c; };

PT<solTSP> algo(1.0f, 60.0f, 12, 300, 3000, 4, 0, 50);
algo.loadCheckpoint("job.ckpt", le);           // continua a execução salva, se o arquivo existir
algo.setCheckpoint("job.ckpt", escreve, 60.0); // salva a cada minuto e ao final
solTSP melhor = algo.start(0, &prob);
\end{lstlisting}

\begin{tabela}{>{\raggedright\arraybackslash}p{0.40\linewidth}>{\raggedright\arraybackslash}X}
\toprule
\cab{Função} & \cab{Efeito} \\
\midrule
\code{setCheckpoint(\allowbreak{}arquivo, escritor, aCadaSegundos)} & Salva o estado ao fim de \code{start()} e, se o intervalo for maior que zero, periodicamente. \\
\code{loadCheckpoint(\allowbreak{}arquivo, leitor)} & O próximo \code{start()} continua a execução salva. Devolve \code{false} se o arquivo não existe e lança \code{std::runtime\_\allowbreak{}error} se ele está danificado. \\
\code{saveCheckpoint(\allowbreak{}arquivo, escritor)} & Salva, sob demanda, o estado do último \code{start()}. \\
\bottomrule
\end{tabela}

Uma execução retomada cumpre o PTL e o tempo restantes. A cópia periódica é tirada por cada réplica ao fim da sua rodada, sem parar as outras; o arquivo é escrito com outro nome e só substitui o anterior quando está completo, de modo que uma queda durante a escrita não destrói o checkpoint anterior.

\begin{lstlisting}[style=saida]
new run:  rounds  912 of 3000   best  11253   time of the job so far   150.3 ms
resumed:  rounds 3000 of 3000   best  11049   time of the job so far   606.7 ms
\end{lstlisting}

\begin{nota}
\textbf{O que a retomada não é.} É uma continuação da busca, não uma repetição bit a bit: os geradores aleatórios e as estatísticas das atualizações dinâmicas recomeçam, e \code{getConvergence()} lista apenas as melhorias encontradas depois da retomada. A cópia periódica pode ter até um intervalo de idade.
\end{nota}

\subsection{Esquema de troca e diagnósticos da escada}

\begin{lstlisting}[style=cpp]
algo.setSwapScheme(PT_SWAP_SEQUENTIAL);   // padrão, como na v1
algo.setSwapScheme(PT_SWAP_EVEN_ODD);     // pares alternados
\end{lstlisting}

\begin{itemize}
\item \textbf{Sequencial}: em cada rodada todos os pares propõem troca, do primeiro ao último.
\item \textbf{Par-ímpar} (\textit{deterministic even-odd}): os pares (0,1), (2,3), $\dots$ propõem em uma rodada e os pares (1,2), (3,4), $\dots$ na seguinte. \end{itemize}

\begin{fonte}
\textbf{Fonte.} O esquema par-ímpar é o de Okabe et al. (2001); Syed et al. (2019) mostram que, aplicado de forma determinística, ele torna o PT não reversível e analisam o ganho em idas e voltas.
\end{fonte}

Dois diagnósticos acompanham qualquer execução:

\begin{itemize}
\item \code{getRoundTrips()}: número de réplicas que foram da primeira temperatura à última e voltaram. É o que leva uma solução encontrada em alta temperatura até as baixas.
\item \code{getCommunicationBarrier()}: soma das taxas de rejeição das trocas, $\sum_i (1 - \mathrm{taxa}_i)$. Mede o quanto o intervalo de temperaturas é difícil de atravessar. \end{itemize}

\begin{fonte}
\textbf{Fonte.} A ida e volta como medida da escada é de Katzgraber et al. (2006); a barreira de comunicação global é definida por Syed et al. (2019).
\end{fonte}

\begin{lstlisting}[style=saida]
replicas  scheme       best   round trips   barrier   exchange rates (min..max)
       6  sequential   7499             1      4.91   0.00..0.03
      12  sequential   7472             4      9.42   0.01..0.32
      24  sequential   7438             5     13.11   0.18..0.64
      48  sequential   7400             4     14.96   0.36..0.81
      48  even-odd     7393             4     14.88   0.40..0.82
\end{lstlisting}

Com 6 réplicas quase toda troca é rejeitada: a barreira está no teto (réplicas $-$ 1). Ela cresce cada vez menos com o número de réplicas (13,1 com 24; 15,0 com 48; 16,3 com 96): quando para de crescer, mais réplicas deixam de ajudar. Nos problemas testados o esquema par-ímpar não foi melhor que o sequencial (\hyperref[s9]{Seção 9}), por isso o padrão continua sendo o da v1.

\subsection{Resultados da execução}

\begin{tabela}{>{\raggedright\arraybackslash}p{0.40\linewidth}>{\raggedright\arraybackslash}X}
\toprule
\cab{Função} & \cab{Devolve} \\
\midrule
\code{getEvaluations()} & Número de vizinhos avaliados. \\
\code{getElapsedMs()}, \code{getTimeToBestMs()} & Tempo total e tempo até a melhor solução. \\
\code{getConvergence()} & Cada nova melhor solução: \code{\{timeMs, ptl, value\}} (a curva de convergência). \\
\code{getMoveAcceptanceRates()} & Taxa de aceitação dos movimentos, uma por temperatura. \\
\code{getSwapAcceptanceRates()} & Taxa de aceitação das trocas, uma por par adjacente. \\
\code{getReplicas()} & Solução final de cada temperatura. \\
\code{getThreadsUsed()} & Threads efetivamente criadas. \\
\code{getMoveStatistics()} & \code{\{proposed, accepted\}} de cada movimento registrado. \\
\code{getRoundTrips()}, \code{getCommunicationBarrier()} & Diagnósticos da Seção 6.7. \\
\code{getStopReason()} & Motivo do término (Seção 6.3). \\
\code{getTiming()} & Tempo de cada componente (Seção 6.9). \\
\bottomrule
\end{tabela}

\begin{lstlisting}[style=saida]
best 7563, found at round 177 of 1000, after 14.1 ms of 74.5 ms
2000000 evaluations (26.9 million/s) with 10 threads, seed 8

temperature           1.00     1.58     2.48     3.91     6.17     9.72    15.33    24.15    38.07    60.00
accepted moves       0.002    0.003    0.003    0.003    0.004    0.005    0.009    0.015    0.027    0.051
final replica         7563     7567     7589     7638     7678     7817     8053     8166     8907    11744
accepted exchanges        0.201    0.115    0.048    0.207    0.133    0.124    0.038    0.017    0.007
\end{lstlisting}

\subsection{Tempo de cada componente}

Depois de \code{setProfiling(true)}, \code{getTiming()} informa para onde foi o tempo da execução: as cadeias de Markov (o código do problema, \code{neighbor}/\code{evaluate}, mais o teste de aceitação), as propostas de troca, as atualizações de temperatura e as threads à espera de um nó pronto. Esses quatro tempos são somados sobre as threads e por isso fecham em threads $\times$ (tempo total $-$ construção). O tempo total e o da construção inicial são sempre medidos. A medição custa duas leituras de relógio por execução de nó, e não por movimento, e vem desligada.

\begin{lstlisting}[style=saida]
long chains (8 replicas, MCL 2000, 8 threads): best 7510 in 62.0 ms
   chains (problem code)    369.39 thread-ms   74.5%
   exchanges                  0.25 thread-ms    0.0%
   temperature updates        0.08 thread-ms    0.0%
   waiting + scheduling     126.10 thread-ms   25.4%

short chains (8 replicas, MCL 20, 8 threads): best 7466 in 202.3 ms
   chains (problem code)    369.42 thread-ms   22.8%
   exchanges                 24.40 thread-ms    1.5%
   temperature updates        1.20 thread-ms    0.1%
   waiting + scheduling    1223.38 thread-ms   75.6%
\end{lstlisting}

As duas execuções fazem o mesmo número de avaliações (o tempo nas cadeias é igual), mas com cadeias curtas as threads se encontram nas trocas cem vezes mais e passam três quartos do tempo esperando umas pelas outras: o MCL é também um parâmetro de eficiência paralela. A leitura prática é esta: se o tempo está nas cadeias, o que vale otimizar é o código do problema (\code{neighborEval}); se está na espera, aumentar o MCL ou usar menos threads.

\subsection{Gerenciamento de experimentos}

O experimento típico de um artigo (R execuções por instância, um arquivo por execução e um resumo com melhor, média, mediana e desvio-padrão) está no cabeçalho opcional \code{Experiment.h}. O usuário escreve apenas a função que faz \textit{uma} execução; a API cuida do laço, das sementes, dos arquivos e das estatísticas.

\begin{lstlisting}[style=cpp]
#include "include/Experiment.h"

PTExperiment exp("resultados");           // pasta de saída
exp.setRuns(30);                          // execuções por instância
exp.setFirstSeed(1);                      // a execução k usa a semente 1 + k
exp.addInstance("instancias/a.txt");      // o que a função abaixo recebe
exp.setReference(1234.0);                 // opcional: melhor valor conhecido ou limite inferior
exp.addInstancesFromFile("lista.txt");    // ou uma linha "caminho [referência]" por instância

exp.run([](const std::string& instancia, unsigned long long semente){
    MeuProblema prob(instancia);
    PT<minhaSolucao> algo(1.0, 60.0, 16, 500, 1000, 4, 0, 50);
    algo.setSeed(semente);
    minhaSolucao melhor = algo.start(0, &prob);
    return ptResult(algo, melhor);        // PTRunResult: tudo o que a API mediu
});
\end{lstlisting}

\begin{tabela}{>{\raggedright\arraybackslash}p{0.26\linewidth}>{\raggedright\arraybackslash}X}
\toprule
\cab{Arquivo} & \cab{Conteúdo} \\
\midrule
\code{<instância>/run\_<k>.txt} & Uma execução: semente, valor, tempos, rodadas, avaliações, motivo do término, escada final, taxas de aceitação, curva de convergência e, se o usuário preencher \code{PTRunResult::solution}, a solução. \\
\code{runs.csv} & Uma linha por execução. \\
\code{summary.csv} & Uma linha por instância: execuções, melhor, média, mediana, desvio-padrão (amostral), pior, tempo médio, tempo médio até a melhor, avaliações e, havendo referência, os desvios percentuais e o número de execuções que a atingiram. \\
\code{summary.txt} & A mesma tabela, alinhada, para ler no terminal. \\
\code{summary.tex} & A mesma tabela em \LaTeX{} (\code{booktabs}). \\
\code{<instância>/\allowbreak{}convergence.csv} & Convergência da instância sobre as suas execuções, por rodada: em cada ponto de uma grade (100 pontos, mais próximos no início; \code{setCurvePoints}), o melhor, a mediana, a média e o pior dos valores que as execuções tinham até ali. \\
\code{<instância>/\allowbreak{}convergence\_time.csv} & O mesmo por tempo, em segundos. \\
\code{<instância>/\allowbreak{}time\_to\_target.csv} & Instâncias com referência: o tempo que cada execução levou para atingi-la, em ordem crescente, com a probabilidade empírica $(i - 0{,}5)/R$ de já tê-la atingido. \\
\code{figures.gp}, \code{figures.tex} & Roteiros que desenham os três arquivos anteriores, com gnuplot ou com \LaTeX{} e pgfplots. \\
\bottomrule
\end{tabela}

Saída de \code{exampleExperiment} (três instâncias, dez execuções cada; as colunas do pior valor e do tempo até a melhor solução foram omitidas aqui). Na primeira instância a referência é um valor-alvo, e não o melhor conhecido; seis das dez execuções o atingiram.

\begin{lstlisting}[style=saida]
instance  runs       best       mean     median       sd  time(s)  gap best%  gap mean%  hits
tsp60       10    5976.00    6020.30    6017.50    36.35     0.01      -0.73       0.00     6
tsp100      10    7435.00    7601.50    7618.50    91.81     0.01       1.16       3.42     0
tsp150      10    9708.00    9837.80    9824.00    82.13     0.01          -          -     -

tsp100/convergence.csv:
   round,runs,best,median,mean,worst
   0,10,31731,33165,33219.2,34412
   1,10,23830,26341,26201.9,27360
tsp60/time_to_target.csv:
   seconds,probability,run
   0.001684818,0.05,9
   0.002932627,0.15,3
tsp60: 6 of 10 executions reached 6020, in 0.0048 s on average
\end{lstlisting}

\textbf{Curvas e tempo até o alvo.} Os três arquivos por instância são CSV simples, feitos para serem desenhados, e a pasta recebe dois roteiros que os desenham: \code{figures.gp} (\code{gnuplot figures.gp} grava \code{convergence.pdf} e \code{time\_to\_target.pdf} em cada instância) e \code{figures.tex} (\code{lualatex figures.tex} grava \code{figures.pdf}, uma figura por página, com pgfplots; cada \code{tikzpicture} pode ser copiado para um artigo). A API não os executa. A figura de convergência traz a mediana, a média e a faixa entre a melhor e a pior execução, com as rodadas em escala logarítmica. As curvas vêm da lista de novas melhores soluções de cada execução (\code{getConvergence}), e por isso os primeiros pontos podem mudar um pouco entre duas execuções com a mesma semente: a ordem em que as réplicas registram as primeiras soluções depende das threads. O arquivo \code{time\_to\_target.csv} contém os dados do gráfico de tempo até o alvo (\textit{time-to-target plot}) de Aiex, Resende e Ribeiro (2007); com a referência usada como alvo, \code{summary.csv} traz também quantas execuções o atingiram e o tempo médio que levaram.

\textbf{Colunas do \code{runs.csv}.} O padrão é \code{instance}, \code{run}, \code{seed}, \code{value}, \code{seconds}, \code{seconds\_to\_best}, \code{rounds} (o PTL em que a execução terminou), \code{best\_round}, \code{evaluations}, \code{threads}, \code{stop} e \code{round\_trips}. A lista pode ser alterada: \code{addCsvColumns} acrescenta colunas e \code{setCsvColumns} substitui a lista inteira, na ordem dada. Além das do padrão, a API sabe preencher \code{path}, \code{evaluations\_per\_second}, \code{improvements}, \code{replicas}, \code{temp\_first}, \code{temp\_last}, \code{temperatures}, \code{temperatures\_initial}, \code{moves\_proposed}, \code{moves\_accepted}, \code{moves\_rate} (um valor por movimento registrado), \code{seconds\_to\_reference}, \code{swap\_rate\_mean}, \code{swap\_rate\_min}, \code{swap\_rate\_max}, \code{barrier}, \code{swap\_rates}, \code{move\_rates}, \code{construction\_seconds}, \code{chains\_ms}, \code{exchanges\_ms}, \code{updates\_ms}, \code{waiting\_ms}, \code{reference}, \code{gap\_percent} e \code{solution} (a lista está em \code{PTExperiment::csvColumnNames()}). Valores do próprio usuário entram com \code{r.add("nome", valor)} dentro da função da execução: cada nome vira mais uma coluna e fica guardado no arquivo da execução, de modo que não se perde em uma retomada.

\begin{lstlisting}[style=cpp]
exp.setCsvColumns({"instance", "seed", "value", "rounds", "stop", "barrier", "gap_percent"});

// dentro da função da execução:
PTRunResult r = ptResult(algo, melhor);
r.add("trocas_de_ferramenta", melhor.trocas);     // uma medida do problema
return r;
\end{lstlisting}

Um experimento interrompido continua de onde parou: as execuções cujo arquivo existe são lidas, e não repetidas (\code{setResume(false)} refaz tudo). Os resumos são regravados ao fim de cada instância, e \code{getSummary()} e \code{getRuns(i)} entregam os mesmos números ao programa. As execuções são feitas uma após a outra, porque cada uma já usa todas as threads que recebe.

\textbf{Várias máquinas.} \code{setShard(indice, partes)} divide o experimento: o processo só faz as execuções da sua parte, de modo que o mesmo programa pode ser iniciado em várias máquinas, gravando na mesma pasta ou em pastas reunidas depois. \code{report()} não executa nada: lê as execuções que estão na pasta e grava todos os relatórios (serve também para refazer o \code{runs.csv} com outras colunas).

\begin{lstlisting}[style=cpp]
exp.setShard(atoi(argv[1]), 4);   // este processo: parte argv[1] de 4
exp.run(execucao);
// depois, uma vez, na pasta com todos os arquivos:
exp.report();
\end{lstlisting}

\subsection{Comparação de configurações e testes estatísticos}

Cada configuração é um experimento sobre as mesmas instâncias. O cabeçalho opcional \code{Comparison.h} os coloca lado a lado com os testes normalmente pedidos em um artigo. A primeira configuração adicionada é o \textit{controle}; as demais são comparadas com ela.

\begin{lstlisting}[style=cpp]
#include "include/Comparison.h"

PTExperiment fixa("res/fixa"), tipo1("res/tipo1");   // run() ou report() em cada um
PTComparison cmp;
cmp.add("fixa", fixa);            // controle
cmp.add("tipo1", tipo1);
cmp.setAlpha(0.05);               // nível das marcas +, -, = (padrão)
cmp.write("res");                 // comparison.txt, .csv, _total.csv e .tex
std::cout << cmp.table();
\end{lstlisting}

\begin{tabela}{>{\raggedright\arraybackslash}p{0.20\linewidth}>{\raggedright\arraybackslash}p{0.34\linewidth}>{\raggedright\arraybackslash}X}
\toprule
\cab{Nível} & \cab{Teste} & \cab{Resultado} \\
\midrule
Cada instância & Mann-Whitney U sobre os valores das execuções das duas configurações & Valor-$p$ e uma marca: \code{+} melhor que o controle, \code{-} pior, \code{=} sem diferença significativa. \\
Todas as instâncias & Wilcoxon de postos sinalizados sobre as médias das instâncias (pareadas por instância) & Valor-$p$ de cada configuração contra o controle; com mais de uma comparação, também corrigido pelo método de Holm. \\
Três ou mais configurações & Friedman sobre os postos das médias em cada instância & Posto médio de cada configuração e um valor-$p$. \\
\bottomrule
\end{tabela}

A tabela traz ainda, por configuração, o desvio médio (em relação à referência das instâncias, se todas tiverem uma; senão, à melhor média da instância), o número de instâncias em que ela tem a melhor média e a contagem de \code{+}, \code{=} e \code{-}. Todos os testes são bilaterais e usam postos, portanto não supõem valores normais. Os valores-$p$ são exatos para amostras pequenas (Mann-Whitney até 60 valores no total, postos sinalizados até 50 instâncias, com empates) e usam a aproximação normal acima disso; Friedman usa a aproximação qui-quadrado. Os testes também são funções independentes: \code{ptstat::mannWhitney}, \code{ptstat::wilcoxonSignedRank}, \code{ptstat::friedman} e \code{ptstat::holm}.

Saída de \code{exampleComparison}: seis instâncias do TSP, dez execuções, 12 réplicas com escada fixa (controle), 12 réplicas com o tipo 1 e 4 réplicas com escada fixa.

\begin{lstlisting}[style=saida]
instance       fixed      rate23         p    four_replicas         p
tsp50        5189.00     5189.00    1.0000 =        5298.30    0.0108 -
tsp70        6302.10     6311.00    0.6955 =        6453.20   <0.0001 -
tsp90        7130.10     7170.60    0.0785 =        7297.50   <0.0001 -
tsp110       7973.90     8025.30    0.3246 =        8136.40    0.0068 -
tsp130       9008.50     9005.00    0.8113 =        9172.50    0.0052 -
tsp150       9837.80     9778.60    0.1596 =        9910.70    0.3423 =

configuration    mean rank     gap %  best mean      +    =    -         p      Holm
fixed                 1.42      0.11          4                    control
rate23                1.58      0.23          3      0    6    0    0.8125    0.8125
four_replicas         3.00      2.02          0      0    1    5    0.0312    0.0625

Friedman test (the configurations have the same mean rank): p = 0.0087
\end{lstlisting}

A leitura é a que se espera: o tipo 1 não se distingue da escada fixa em nenhuma instância, e quatro réplicas são piores em cinco das seis. O exemplo mostra também um limite do método: com seis instâncias, o menor valor-$p$ possível do teste de postos sinalizados é $2/2^6 = 0{,}031$, que a correção de Holm para duas comparações leva a 0,0625; com menos de seis instâncias esse teste não chega a 0,05, qualquer que seja o resultado.

\begin{fonte}
\textbf{Fonte.} Testes de Mann e Whitney (1947), Wilcoxon (1945) e Friedman (1937) e correção de Holm (1979). A metodologia (postos sinalizados sobre os conjuntos de dados, Friedman e postos médios para vários algoritmos, em vez de testes paramétricos) é a recomendada por Demšar (2006). A implementação é própria e foi conferida com o SciPy e com a enumeração de todos os casos em amostras pequenas (diferenças abaixo de $10^{-14}$).
\end{fonte}

\subsection{Linha de comando e calibração com o irace}

O cabeçalho opcional \code{Options.h} lê todos os parâmetros do PT da linha de comando, no formato \code{-{}-NOME valor}, com os nomes dos exemplos da v1 (\code{-{}-TEMP\_INIT}, \code{-{}-TEMP\_FIM}, \code{-{}-N\_REPLICAS}, \code{-{}-MCL}, \code{-{}-PTL}, \code{-{}-TEMP\_DIST}, \code{-{}-TYPE\_UPDATE}, \code{-{}-TEMP\_UPDATE}, \code{-{}-THREAD\_USED}) e os da v2 (\code{-{}-SEED}, \code{-{}-TIME\_LIMIT}, \code{-{}-TARGET}, \code{-{}-STAGNATION\_ROUNDS}, \code{-{}-STAGNATION\_TIME}, \code{-{}-SWAP\_SCHEME}, \code{-{}-CALIBRATE}, entre outros). As opções que a API não conhece ficam guardadas para o problema, de modo que os parâmetros do problema são calibrados do mesmo jeito.

\begin{lstlisting}[style=cpp]
#include "include/Options.h"

int main(int argc, char* argv[]){
    PTOptions opt(argc, argv);                       // opt.instance = primeiro argumento
    MeuProblema prob(opt.instance, opt.get("--ALFA", 0.5));
    minhaSolucao melhor = opt.run<minhaSolucao>(&prob);   // monta o PT, calibra se pedido, executa
    if(!opt.printForTuner()) std::cout << melhor.evalSol << "\n";
}
\end{lstlisting}

Com \code{-{}-IRACE 1} o programa imprime apenas o custo da execução, que é o que o irace (López-Ibáñez et al., 2016) lê; \code{-{}-IRACE 2} imprime custo e tempo. A opção \code{-{}-IRACE\_COST} define o que é minimizado, o que permite calibrar pela função objetivo, pelo tempo de execução ou pelos dois:

\begin{tabela}{l>{\raggedright\arraybackslash}X}
\toprule
\cab{\code{-{}-IRACE\_COST}} & \cab{Custo informado ao irace} \\
\midrule
\code{value} (padrão) & O valor da função objetivo ao fim da execução. Todas as configurações recebem o mesmo \code{-{}-TIME\_LIMIT}; a qualidade é comparada sob o mesmo tempo. \\
\code{time} & Os segundos até atingir \code{-{}-TARGET}. Quem não atinge o alvo dentro do limite paga \code{-{}-IRACE\_PENALTY} (padrão 10) vezes o limite. \\
\code{value+time} & valor $+\ w\ \times$ segundos, com $w$ dado por \code{-{}-IRACE\_TIME\_WEIGHT}: um preço por segundo, para execuções que param sozinhas (estagnação ou PTL). \\
\bottomrule
\end{tabela}

A pasta \code{irace/} traz um cenário completo para o exemplo SSP: \code{parameters.txt} (os parâmetros e as faixas), \code{scenario.txt}, \code{target-runner} (o script que o irace chama, onde fica o orçamento de cada execução) e \code{ptSSP.cpp} (o programa, em que só duas linhas dependem do problema). \code{make check} chama o \code{target-runner} uma vez, como o irace faz, e \code{make tune} executa o irace. O cenário foi executado de ponta a ponta com o irace 4.5.1: 300 execuções de 1 s nas dez instâncias de SSP, em 165 s, com este resultado:

\begin{lstlisting}[style=saida]
# Best configurations as commandlines (listed from best to worst according to the sum of ranks):
25 --TEMP_INIT 0.312 --TEMP_FIM 2.828 --N_REPLICAS 7 --MCL 101 --TEMP_DIST 4 --TYPE_UPDATE 5 --TEMP_UPDATE 0 --SWAP_SCHEME 0
16 --TEMP_INIT 0.155 --TEMP_FIM 1.062 --N_REPLICAS 12 --MCL 1150 --TEMP_DIST 3 --TYPE_UPDATE 5 --TEMP_UPDATE 50 --SWAP_SCHEME 1
26 --TEMP_INIT 0.168 --TEMP_FIM 1.033 --N_REPLICAS 22 --MCL 122 --TEMP_DIST 4 --TYPE_UPDATE 1 --TEMP_UPDATE 20 --SWAP_SCHEME 1
\end{lstlisting}

Esse orçamento serviu para validar a integração, não para recomendar parâmetros: 300 execuções são poucas para oito parâmetros.

\section{Atualização dinâmica das temperaturas: os sete tipos}\label{s7}

O sétimo argumento do construtor escolhe a regra e o oitavo, o intervalo $\upsilon$ (propostas de troca entre duas atualizações). A v1 tinha os tipos 0 a 3; a v2 mantém os quatro e acrescenta três.

\begin{tabela}{r>{\raggedright\arraybackslash}X>{\raggedright\arraybackslash}Xll}
\toprule
\cab{Tipo} & \cab{Regra} & \cab{Extremos da escada} & \cab{Reprodutível} & \cab{Situação} \\
\midrule
0 & Temperaturas fixas & fixos & sim & v1 \\
1 & Taxa de troca de 23\% & fixos & sim ($\upsilon \ge 2$) & v1, regra revista \\
2 & Taxas de troca iguais & fixos & sim & v1, com proteção \\
3 & \textit{Feedback-optimized}, como na tese & fixos (com 3 ou mais réplicas) & não & v1, inalterada \\
4 & \textit{Feedback-optimized} por quantis & fixos & não & nova \\
5 & Aproximação estocástica para 23,4\% & primeira fixa; última $\le$ máxima inicial & sim & nova \\
6 & Taxas de rejeição iguais & fixos & não & nova \\
\bottomrule
\end{tabela}

Cada par mede, na janela de $\upsilon$ propostas, a sua taxa de troca e a energia média das duas réplicas ($\bar{E}$). Os tipos 3 e 4 usam o fluxo de réplicas rotuladas; os tipos 5 e 6 usam a probabilidade de troca calculada em cada proposta, e não a contagem de aceites, que é uma estimativa mais ruidosa da mesma grandeza.

\subsection*{Tipo 1 — taxa de troca de 23\%}

Para cada par, com $\Delta \bar{E}$ a diferença das energias médias e $\Delta \beta$ a distância entre as temperaturas em 1/$T$, a distância que daria 23\% de trocas é obtida de

\[
0{,}23 = \exp(-\Delta\beta^{*}\,\Delta\bar{E}) \quad\Longrightarrow\quad \Delta\beta^{*} = \frac{-\ln 0{,}23}{\Delta\bar{E}}.
\]

Se a taxa medida está entre 20\% e 26\%, nada é feito. Fora dessa faixa, o par move a sua distância na direção de $\Delta \beta$* (mais estreita abaixo de 20\%, mais larga acima de 26\%), aplicando metade da correção. A correção é dividida entre as duas temperaturas do par; a primeira e a última temperatura da escada nunca se movem, e nenhuma temperatura ultrapassa o ponto médio até a vizinha.

\begin{fonte}
\textbf{Fonte.} O alvo de 23,4\% é de Atchadé, Roberts e Rosenthal (2011), na linha de Roberts, Gelman e Gilks (1997); a equação $0{,}23 = e^{\Delta\beta\,\Delta\bar{E}}$ é a da tese. A ancoragem dos extremos, a correção nos dois sentidos, a faixa morta e o amortecimento são adaptações desta versão: a regra da v1 só agia abaixo de 20\% e só aproximava temperaturas, o que fazia a escada inteira descer.
\end{fonte}

\subsection*{Tipo 2 — taxas de troca iguais}

A temperatura $T_{i}$ é recolocada entre as vizinhas de modo que a probabilidade de troca com a de baixo seja igual à probabilidade de troca com a de cima, $P(r_{i-1} \to r_i) = P(r_i \to r_{i+1})$, usando as energias médias da janela:

\[
T_i = \frac{\Delta\bar{E}_{i-1} + \Delta\bar{E}_i}{\dfrac{\Delta\bar{E}_{i-1}}{T_{i-1}} + \dfrac{\Delta\bar{E}_i}{T_{i+1}}}.
\]

Na v2 o novo valor só é aceito se ficar entre as duas vizinhas. Sem essa proteção, com estimativas ruidosas de $\Delta \bar{E}$, a regra produzia escadas fora de ordem e fora do intervalo (em um teste com intervalo de 1 a 60 surgiram temperaturas de 0,33 e 65).

\begin{fonte}
\textbf{Fonte.} Hukushima (1999), como implementado na tese. A proteção de ordem é desta versão.
\end{fonte}

\subsection*{Tipo 3 — \textit{feedback-optimized}, como na tese}

Cada temperatura registra a fração $f(T)$ das réplicas rotuladas que a visitaram vindo da primeira temperatura («subindo»). O método procura a escada que maximiza as idas e voltas, tornando a densidade de temperaturas proporcional a

\[
\eta(T) \propto \sqrt{\frac{\mathrm{d}f/\mathrm{d}T}{\Delta T}}.
\]

A implementação da tese percorre a escada da primeira para a última temperatura, definindo cada temperatura a partir da anterior com a constante de normalização calculada pelo par da primeira temperatura. Foi mantida exatamente como estava; a única alteração é que os valores lidos por outros pares passaram a ser variáveis atômicas.

\begin{fonte}
\textbf{Fonte.} Katzgraber, Trebst, Huse e Troyer (2006), na forma descrita na tese.
\end{fonte}

\subsection*{Tipo 4 — \textit{feedback-optimized} por quantis}

Usa a mesma densidade do tipo 3, mas recalcula a escada inteira de uma vez, com os extremos fixos. A integral da densidade em cada intervalo [$T_{k}$, $T_{k+1}$] é

\[
w_k = \sqrt{\,|f_k - f_{k+1}|\,},
\]

e as novas temperaturas são os pontos que dividem a soma acumulada dos $w_{k}$ em $M-1$ partes iguais, com interpolação linear em $T$ dentro de cada intervalo. A atualização só ocorre quando todas as temperaturas têm dados de fluxo, e aplica metade do passo (a nova temperatura é a média entre a atual e a calculada), o que mantém a escada ordenada.

\begin{fonte}
\textbf{Fonte.} A densidade ótima é a de Katzgraber et al. (2006). A formulação por quantis com extremos fixos e meio passo é desta versão.
\end{fonte}

\subsection*{Tipo 5 — aproximação estocástica}

A escada é parametrizada pelas razões entre temperaturas vizinhas:

\[
T_{i+1} = T_i \exp\bigl(\exp(\rho_i)\bigr).
\]

A cada atualização, o par $i$ corrige o seu parâmetro com um passo decrescente, em que $\bar{p}_{i}$ é a média das probabilidades de troca do par na janela e $n$ conta as atualizações já feitas:

\[
\rho_i \leftarrow \rho_i + (n+1)^{-0{,}6}\,(\bar{p}_i - 0{,}234).
\]

Se o par troca mais que 23,4\%, $\rho$ cresce e as temperaturas se afastam; se troca menos, se aproximam. Como o passo diminui, a escada converge em vez de oscilar. A primeira temperatura é fixa e todas as outras decorrem dos $\rho$; a dupla exponencial garante $T_{i+1}$ > $T_{i}$. Os $\rho$ iniciais vêm da escada inicial, que por isso precisa ser crescente (distribuições 1, 3 e 4).

Cada par só escreve a sua temperatura mais quente, depois que o par anterior executou; por isso este é o único dos tipos novos reprodutível com semente.

\begin{fonte}
\textbf{Fonte.} Miasojedow, Moulines e Vihola (2013): a parametrização, a regra de Robbins–Monro e o alvo de 23,4\% são do artigo, que também prova a convergência da adaptação. Duas escolhas são desta versão: o expoente 0,6 do passo (o artigo exige um expoente entre 0,5 e 1) e o limite da última temperatura à máxima da escada inicial, sem o qual ela dispara em problemas com réplicas de sobra. O artigo também adapta a distribuição de proposta dos movimentos, o que não foi implementado, por depender do problema.
\end{fonte}

\subsection*{Tipo 6 — taxas de rejeição iguais}

Seja $r_i = 1 - \bar{p}_i$ a rejeição média do par $i$ na janela. A soma acumulada

\[
\Lambda_k = r_0 + r_1 + \dots + r_{k-1}
\]

é tratada como função de $\beta = 1/T$, com interpolação linear entre as temperaturas. As novas temperaturas são os pontos em que essa função vale $k/(M-1)$ do total: cada par passa a rejeitar a mesma fração das trocas. Os extremos ficam fixos e é aplicada metade do passo, em $\beta$. A atualização é feita pelo par da primeira temperatura para a escada inteira.

\begin{fonte}
\textbf{Fonte.} Syed, Bouchard-Côté, Deligiannidis e Doucet (2019): a escada ótima é a que iguala as rejeições, e o procedimento de ajuste é a equalização da barreira de comunicação acumulada. O artigo usa interpolação monotônica e refaz a escada entre rodadas de ajuste; a interpolação linear em 1/$T$, o meio passo e a atualização durante a própria execução são simplificações desta versão.
\end{fonte}

\subsection*{Efeito de cada regra sobre uma escada ruim}

O exemplo \code{exampleTempUpdate} parte de uma escada linear de 1 a 60 com 10 réplicas (quase sem trocas no extremo frio) e imprime o que cada regra faz (TSP de 100 cidades, 2000 rodadas, $\upsilon = 20$). As taxas são as da execução inteira, incluindo as rodadas anteriores à estabilização.

\begin{lstlisting}[style=saida]
inicial            1.00   7.56  14.11  20.67  27.22  33.78  40.33  46.89  53.44  60.00
tipo 0  trocas     0.011  0.022  0.068  0.158  0.178  0.214  0.336  0.306  0.376

tipo 1  final      1.00   1.23   2.03   3.64   6.38   9.84  15.09  24.62  43.66  60.00
        trocas     0.241  0.242  0.247  0.222  0.194  0.113  0.054  0.035  0.036
tipo 2  final      1.00   1.30   1.99   2.59  11.24  14.90  25.28  34.72  41.17  60.00
        trocas     0.095  0.068  0.107  0.146  0.194  0.239  0.258  0.313  0.297
tipo 3  final      1.00   2.67  17.23  21.88  23.73  27.17  27.17  40.34  40.34  60.00
        trocas     0.006  0.008  0.214  0.677  0.416  0.909  0.103  0.902  0.040
tipo 4  final      1.00  17.59  23.96  30.62  37.89  43.10  46.89  50.60  54.88  60.00
        trocas     0.007  0.154  0.200  0.242  0.420  0.571  0.600  0.514  0.512
tipo 5  final      1.00   1.16   1.46   1.91   2.34   3.21   5.32   6.93   9.89  25.43
        trocas     0.013  0.146  0.232  0.182  0.211  0.284  0.198  0.161  0.200
tipo 6  final      1.00   1.15   1.94   3.99   7.92  13.22  21.62  31.91  45.60  60.00
        trocas     0.052  0.084  0.099  0.078  0.055  0.085  0.071  0.080  0.064
\end{lstlisting}

Leitura: o tipo 6 iguala as taxas (5\% a 10\%) mantendo os dois extremos, e mostra que 10 réplicas são poucas para ir de 1 a 60 com boas trocas; o tipo 5 busca 23\% e, para isso, encolhe a escada (termina em 25,4); o tipo 1 chega a 23\% no lado frio e deixa o desvio no lado quente; o tipo 4, que otimiza o fluxo e não a taxa, concentra as temperaturas no lado quente nesta instância.

\subsection*{Intervalo automático entre as atualizações}

O oitavo argumento do construtor é o número de propostas de troca entre duas atualizações. Com \code{PT\_UPDATE\_AUTO} (o valor 0) a API escolhe: a primeira atualização ocorre depois de 10 propostas e cada intervalo é o dobro do anterior (10, 20, 40, \dots). A escada é corrigida muitas vezes no começo, quando está longe de uma boa, e depois com estatísticas de janelas mais longas; e há um parâmetro a menos para calibrar. \code{setAutoUpdate(primeiro, fator)} muda os dois números. Vale para os seis tipos e mantém a reprodutibilidade dos tipos 1, 2 e 5. Na v1 o valor 0 causava uma divisão por zero, de modo que nenhum código existente muda de comportamento.

\begin{fonte}
\textbf{Fonte.} As rodadas de ajuste de comprimento dobrado são as do procedimento de ajuste de Syed et al. (2019). Aplicá-las às seis regras da API é desta versão. O resultado medido está na \hyperref[s9]{Seção 9.5}: equivalente ao intervalo fixo, sem ganho de qualidade.
\end{fonte}

\section{Calibração automática das temperaturas}\label{s8}

As duas funções desta seção escolhem temperaturas antes da execução, a partir de medidas feitas no próprio problema. São chamadas entre a construção do objeto e o \code{start()}, e usam a semente de \code{setSeed()}.

\subsection{\code{calibrateTemperatures}: primeira e última temperatura}

\begin{lstlisting}[style=cpp]
std::pair<double,double> t = algo.calibrateTemperatures(&prob);                // padrões
std::pair<double,double> t = algo.calibrateTemperatures(&prob, 20000, 0.10, 0.0003);
//                                             (problema, amostras, aceitaQuente, aceitaFria)
\end{lstlisting}

Funcionamento:

\begin{enumerate}
\item A partir de uma solução construída, faz uma descida de \code{amostras} movimentos (padrão 5000), aceitando só os que não pioram.
\item Guarda os aumentos de custo $\delta$ > 0 dos movimentos rejeitados na segunda metade da descida, quando a solução já está perto de um ótimo local. São os movimentos que uma réplica fria de fato encontra.
\item Define $A(T)$ como a média de $\exp(-\delta/T)$ sobre essa amostra: a fração dos movimentos de piora que uma réplica na temperatura $T$ aceitaria.
\item Resolve, por bisseção em log $T$, $A(T_{\mathrm{min}}) = {}$\code{aceitaFria} (padrão 0,001) e $A(T_{\mathrm{max}}) = {}$\code{aceitaQuente} (padrão 0,02), e reconstrói a escada com a distribuição e o número de réplicas do construtor. \end{enumerate}

O custo são \code{amostras} avaliações, não contadas em \code{getEvaluations()}. Para um alvo de 23\% na aceitação dos movimentos da réplica mais quente, usa-se \code{calibrateTemperatures(\allowbreak{}\&prob, 5000, 0.23)}.

\begin{fonte}
\textbf{Fonte.} Procedimento próprio desta versão. A ideia de obter a temperatura que produz uma taxa de aceitação desejada a partir de uma amostra de movimentos de piora é a mesma usada para a temperatura inicial do \textit{simulated annealing} (Ben-Ameur, 2004); a amostragem na segunda metade de uma descida e os alvos padrão foram escolhidos aqui, por experimento nos problemas TSP, SSP e IPMTC.
\end{fonte}

\begin{lstlisting}[style=saida]
defaults                       Tmin   2.706  Tmax   41.214   best 7467
acceptHot 0.10, cold 0.0003    Tmin   1.847  Tmax  239.485   best 7503
acceptHot 0.23                 Tmin   2.706  Tmax  427.225   best 7424
by hand                        Tmin   1.000  Tmax   60.000   best 7485
\end{lstlisting}

Essas linhas são uma execução de cada opção; a comparação com 20 sementes está na \hyperref[s9]{Seção 9}.

\subsection{\code{calibrateLadder}: escada para uma taxa de troca}

\begin{lstlisting}[style=cpp]
std::vector<double> escada = algo.calibrateLadder(&prob);                       // 23%
std::vector<double> escada = algo.calibrateLadder(&prob, 0, 0.23, 1000, 30, 2.0);
//                               (problema, threads, taxa, rodadas, iteracoes, Tmin)
\end{lstlisting}

Mantém a primeira temperatura e procura todas as outras com execuções-piloto do próprio PT. A escada é escrita como $T_{i+1} = T_i\,e^{e^{\rho_i}}$ e os pilotos usam a aproximação estocástica do tipo 5 com a taxa pedida como alvo:

\begin{enumerate}
\item A cada 10 propostas de troca, cada par corrige o seu $\rho_i$ em $\gamma_n(\bar{p}_i - \mathit{alvo})$, em que $\bar{p}_i$ é a probabilidade média de troca do par desde a correção anterior e $\gamma_n = (n+1)^{-0{,}6}$: a distância cresce se o par troca acima do alvo e diminui se troca abaixo.
\item São feitos \code{iteracoes} pilotos (padrão 10) de \code{rodadas} rodadas (padrão 200), cada um continuando das réplicas e da escada do anterior; o passo $\gamma_n$ continua decrescendo de um piloto para o outro.
\item A escada devolvida usa a média dos $\rho_i$ na segunda metade dos pilotos, o que filtra o ruído da medida. \end{enumerate}

A última temperatura é um resultado: é onde o número de réplicas chega com aquela taxa. Se há mais réplicas do que a taxa pede, a escada para na temperatura que já aceita 90\% dos movimentos (um passeio aleatório, acima do qual temperaturas mais altas nada acrescentam) e as réplicas restantes ficam próximas dela. Sem esse limite a escada cresceu sem controle nos problemas SSP e IPMTC (temperatura máxima de $10^{65}$). Os pilotos custam iterações $\times$ rodadas $\times$ réplicas $\times$ MCL avaliações e ignoram os critérios de parada, o \textit{callback} e o checkpoint do objeto.

As duas calibrações se combinam: \code{calibrateTemperatures} define a primeira temperatura e \code{calibrateLadder}, as demais.

\begin{fonte}
\textbf{Fonte.} A regra de correção é a aproximação estocástica de Miasojedow, Moulines e Vihola (2013), a mesma do tipo 5; o alvo padrão de 23\% vem de Atchadé, Roberts e Rosenthal (2011). São próprios desta versão o uso em pilotos separados da execução, a média dos $\rho_i$ na segunda metade e o limite de 90\% de aceitação.
\end{fonte}

\begin{lstlisting}[style=saida]
geometric ladder from 1 to 60: best 7592, round trips 1, mean exchange rate 0.11
   exchange rate    0.064   0.127   0.175   0.228   0.186   0.115   0.051   0.021   0.003
calibrateLadder, 23%: best 7597, round trips 3, mean exchange rate 0.20
   ladder            1.00    1.18    1.36    1.98    2.67    4.64    6.73    9.59   11.22   15.13
   exchange rate    0.023   0.165   0.054   0.088   0.157   0.289   0.276   0.555   0.161
calibrateLadder, 23%, 10 pilot runs of 200 rounds, from T = 20: best 7621, round trips 10, mean exchange rate 0.24
   ladder           20.00   24.47   30.33   36.23   43.07   50.03   58.08   67.41   77.55   90.62
   exchange rate    0.286   0.179   0.207   0.273   0.288   0.262   0.229   0.265   0.198
   largest deviation of a pair from 23%: 0.058
\end{lstlisting}

\begin{nota}
\textbf{Precisão medida} (TSP de 100 cidades, 10 réplicas, 10 sementes, alvo de 23\%, taxas medidas em uma execução de 5000 rodadas depois da calibração). A partir de $T = 20$, onde as réplicas continuam se movendo, o desvio médio de um par em relação ao alvo foi de 0,025 com os pilotos padrão e de 0,021 com 30 pilotos de 1000 rodadas; com a escada fixa a taxa de um par já varia 0,02 entre sementes, de modo que o resultado está perto do que a medida permite. A regra anterior desta função (correção inteira a cada piloto, pelo fator $\sqrt{\ln(\mathit{alvo})/\ln(\mathit{taxa})}$) dava 0,052 nas mesmas condições. A partir de $T = 1$ as réplicas mais frias ficam presas em ótimos locais: a taxa de um par passa a depender das duas soluções que ele guarda e varia 0,10 entre sementes com a mesma escada, e nenhuma escada cumpre o alvo em toda execução (desvio médio de 0,09; é o caso da segunda saída acima). Ver \hyperref[s11]{Seção 11}.
\end{nota}

\section{Resultados medidos}\label{s9}

Todos os valores são a média da melhor solução em 20 sementes (menor é melhor), com o desvio-padrão entre parênteses quando foi registrado. As diferenças menores que um desvio-padrão não devem ser lidas como ganho ou perda.

\subsection{Tipos de atualização no TSP}

TSP euclidiano aleatório, 16 réplicas, escada geométrica de 1 a 60, PTL 1000; MCL 200 para 100 cidades e 500 para 200 cidades. Os desvios-padrão ficaram entre 26 e 99.

\begin{tabela}{r>{\raggedright\arraybackslash}Xrr}
\toprule
\cab{Tipo} & \cab{Regra} & \cab{TSP 100} & \cab{TSP 200} \\
\midrule
0 & fixa & 7480,9 & 11134,7 \\
1 & 23\% & 7469,9 & 11119,5 \\
2 & taxas iguais & 7504,2 & 11128,0 \\
3 & \textit{feedback} (tese) & 7510,8 & 11188,0 \\
4 & \textit{feedback} (quantis) & 7460,6 & 11113,8 \\
5 & aproximação estocástica & 7478,1 & 11078,6 \\
6 & rejeição igual & 7478,9 & 11128,5 \\
\bottomrule
\end{tabela}

Quase tudo está dentro do ruído. O único resultado que se destaca é o tipo 5 no TSP de 200 cidades (11078,6 com desvio 46, contra 11134,7 da escada fixa). A proteção de ordem do tipo 2 melhorou o seu resultado no TSP 100 de 7562 para 7504.

\subsection{Tipos 5 e 6 nos exemplos da v1}

16 réplicas, MCL 400, PTL 400, intervalo de atualização 20; SSP com temperaturas de 0,2 a 1.

\begin{tabela}{>{\raggedright\arraybackslash}Xrrr}
\toprule
\cab{Problema (instância)} & \cab{Tipo 0 (fixa)} & \cab{Tipo 5} & \cab{Tipo 6} \\
\midrule
SSP (F1001) & 251,55 & 252,10 & 251,50 \\
IPMTC (Large1) & 2189,70 & 2196,20 & 2201,35 \\
\bottomrule
\end{tabela}

Nenhum dos dois superou a escada fixa escolhida à mão nesses problemas; as diferenças estão dentro do ruído (desvio de 20 a 27 no IPMTC).

\subsection{Calibração}

\begin{tabela}{>{\raggedright\arraybackslash}Xrrr}
\toprule
\cab{Problema} & \cab{Temperaturas à mão} & \cab{\code{calibrateTemperatures}} & \cab{\code{calibrateLadder}} \\
\midrule
TSP 100 & 7480,9 (26,4) & 7456,0 (30,1) & 7457,3 (24,6) \\
SSP F1001 & 251,55 (1,16) & 251,75 (1,04) & 251,50 \\
IPMTC Large1 & 2189,7 (27,0) & 2201,3 (28,9) & 2198,9 (27,6) \\
\bottomrule
\end{tabela}

A calibração chega a um resultado equivalente ao das temperaturas escolhidas à mão, sem exigir conhecimento do problema; não as supera de forma significativa. Os valores de \code{calibrateLadder} desta tabela foram medidos com a regra de correção anterior à atual e não foram repetidos; o do IPMTC é também anterior ao limite de 90\% de aceitação.

\subsection{Esquema de troca}

\begin{tabela}{>{\raggedright\arraybackslash}Xrr}
\toprule
\cab{Problema} & \cab{Sequencial} & \cab{Par-ímpar} \\
\midrule
TSP 100 & 7480,9 & 7476,9 \\
TSP 200 & 11134,7 & 11110,6 \\
SSP F1001 & 251,55 & 251,40 \\
IPMTC Large1 & 2189,7 & 2192,0 \\
\bottomrule
\end{tabela}

Não houve diferença significativa, e o número de idas e voltas foi igual ou menor com o esquema par-ímpar. O ganho descrito por Syed et al. (2019) é para amostragem em equilíbrio com muitas réplicas; nestes testes de otimização, com poucas rodadas, ele não apareceu.

\subsection{Intervalo de atualização automático}

Mesma configuração da Seção 9.1; o intervalo fixo é de 50 propostas. Desvio-padrão entre parênteses.

\begin{tabela}{rlrrrr}
\toprule
\cab{Tipo} & \cab{Regra} & \cab{TSP 100, fixo} & \cab{TSP 100, automático} & \cab{TSP 200, fixo} & \cab{TSP 200, automático} \\
\midrule
1 & 23\% & 7470,7 (45) & 7481,8 (30) & 11118,8 (78) & 11109,8 (96) \\
2 & taxas iguais & 7506,8 (43) & 7521,5 (59) & 11123,5 (103) & 11176,6 (75) \\
3 & \textit{feedback} (tese) & 7519,2 (74) & 7508,8 (47) & 11182,5 (117) & 11178,2 (71) \\
4 & \textit{feedback} (quantis) & 7465,3 (47) & 7474,6 (42) & 11120,1 (119) & 11155,6 (130) \\
5 & aproximação estocástica & 7473,8 (40) & 7504,9 (30) & 11084,5 (50) & 11106,1 (114) \\
6 & rejeição igual & 7506,1 (45) & 7499,4 (40) & 11114,1 (74) & 11159,3 (90) \\
\bottomrule
\end{tabela}

O intervalo automático ficou equivalente ao fixo dentro do ruído, com tendência a ser um pouco pior (na média dos seis tipos, $+8$ com 100 cidades e $+24$ com 200); a diferença mais nítida é a do tipo 5 com 100 cidades. Ele é uma conveniência, que elimina um parâmetro, e não um meio de obter soluções melhores; por isso não é o padrão. Esta tabela foi medida com a versão atual e repete, para o intervalo fixo, a medição da Seção 9.1; as pequenas diferenças entre as duas (por exemplo, 7506,1 contra 7478,9 no tipo 6) vêm dos tipos não reprodutíveis e dão a dimensão do ruído.

\section{Testes, exemplos e estrutura de pastas}\label{s10}

\begin{lstlisting}[style=saida]
PTAPI/
 ├─ include/            núcleo da API (só cabeçalhos)
 │    PT.h              classe PT: configuração, start(), calibrações, checkpoint, resultados
 │    Problem.h         interface do problema, geradores aleatórios, movimentos
 │    NodeMCMC.h        nó MC          NodeSwap.h   nó SW          NodeTempUp.h   nó TU
 │    Consumer.h        estado da execução (melhor solução, parada)     Pool.h   threads
 │    Experiment.h      gerenciamento de experimentos (opcional)
 │    Comparison.h      comparação de experimentos e testes estatísticos (opcional)
 │    Options.h         parâmetros pela linha de comando, para o irace (opcional)
 ├─ example/            quickstart.cpp, exampleTSP.cpp e SSP, IPMTC, RCPMS (exemplos da v1, no estilo da v2)
 │   └─ functions/      18 exemplos, um por grupo de funções (make run)
 ├─ instances/          instâncias dos exemplos
 ├─ tests/              testes unitários (make test | tsan | asan | all)
 ├─ irace/              calibração dos parâmetros com o irace (cenário, target-runner, exemplo)
 ├─ bench/              comparação de tempo com a v1 (make bench | bench-sd)
 ├─ v1/                 núcleo da v1, mantido para a comparação de tempo
 └─ docs/               manual (Markdown, em inglês) e report/: este documento e os arquivos que o geram
\end{lstlisting}

\subsection*{Testes}

A suíte \code{tests/\allowbreak{}test\_\allowbreak{}ptapi.cpp} tem 44 testes, sobre três problemas pequenos com codificações diferentes (permutação, binária e contínua). Cobre: as quatro distribuições de temperatura; todos os critérios de parada; a reprodutibilidade com semente para qualquer número de threads; a equivalência das formas de \code{evaluate}; os seis tipos de atualização dinâmica, com intervalo fixo e automático, e a ancoragem dos extremos; as duas calibrações; movimentos com probabilidades; checkpoint e retomada, incluindo arquivos danificados; os dois esquemas de troca; os critérios de parada novos e o motivo do término; a construção por réplica, também em paralelo; o tempo por componente; o gerenciador de experimentos, incluindo a retomada; as opções de linha de comando; as estatísticas; o tratamento de erros; e uma verificação estatística de que cada temperatura amostra a distribuição de Boltzmann, nos dois esquemas de troca. Deve ser executada depois de qualquer mudança no núcleo (\code{make test}) e, antes de uma liberação, também com \code{make tsan} e \code{make asan}.

\subsection*{Exemplos por função}

Em \code{example/\allowbreak{}functions/\allowbreak{}}, \code{make run} compila e executa todos; cada programa devolve 0 somente se o comportamento descrito no seu cabeçalho se confirmou.

\begin{tabela}{>{\raggedright\arraybackslash}p{0.34\linewidth}>{\raggedright\arraybackslash}X}
\toprule
\cab{Programa} & \cab{Funções} \\
\midrule
\code{exampleSeed} & \code{setSeed}, \code{getSeed}, \code{ptapi::rng}, \code{randInt}, \code{rand01} \\
\code{exampleStop} & \code{setTimeLimit}, \code{setTargetValue}, \code{setPTLStopEnd}, \code{setOnImprove}, \code{requestStop} \\
\code{exampleEvaluate} & as três formas de \code{evaluate} e \code{neighborEval} \\
\code{exampleTemperatures} & as quatro distribuições, \code{setTemperatures}, uma única réplica \\
\code{exampleInitialSolutions} & \code{setInitialSolutions}, \code{getReplicas} \\
\code{exampleTempUpdate} & os sete tipos de atualização \\
\code{exampleCalibrateTemperatures} & \code{calibrateTemperatures} \\
\code{exampleCalibrateLadder} & \code{calibrateLadder}, sozinha e combinada \\
\code{exampleMoves} & \code{setMoveProbabilities}, \code{chooseMove}, \code{getMoveStatistics} \\
\code{exampleCheckpoint} & \code{setCheckpoint}, \code{loadCheckpoint}, \code{saveCheckpoint} \\
\code{exampleSwapScheme} & \code{setSwapScheme}, \code{getRoundTrips}, \code{getCommunicationBarrier} \\
\code{exampleStatistics} & \code{getConvergence}, avaliações, tempos, taxas de aceitação \\
\code{exampleV1Compatibility} & classe escrita para a v1, \code{setV1Compatibility} \\
\code{exampleAutoUpdate} & \code{PT\_UPDATE\_AUTO}, \code{setAutoUpdate} \\
\code{exampleTiming} & \code{setProfiling}, \code{getTiming} \\
\code{exampleExperiment} & \code{PTExperiment}, \code{PTRunResult}, \code{ptResult}, curvas, \code{setShard}, \code{report} \\
\code{exampleComparison} & \code{PTComparison} e os testes de \code{ptstat} \\
\code{exampleOptions} & \code{PTOptions} \\
\bottomrule
\end{tabela}

\subsection*{Exemplos da v1}

As cópias de SSP, IPMTC e RCPMS em \code{example/\allowbreak{}} foram reescritas no estilo da v2: \code{ptapi::rng()} em vez de um gerador por chamada, \code{double evaluate(\allowbreak{}const S\&)}, sem \code{Nup}/\code{Ndown} e sem \code{omp.h}, com o movimento aplicado ao próprio argumento de \code{neighbor()}. Os programas principais ganharam as opções \code{--SEED}, \code{--TIME\_\allowbreak{}LIMIT}, \code{--CALIBRATE 1|2|3} e \code{--SWAP\_\allowbreak{}SCHEME 0|1}, e imprimem as estatísticas da execução. Os exemplos originais, com \code{double evaluate(\allowbreak{}S)}, continuam compilando sem alteração.

\section{Limitações conhecidas e pontos em aberto}\label{s11}

\begin{itemize}
\item \textbf{\code{calibrateLadder} só cumpre a taxa pedida onde as réplicas se movem.} Em temperaturas em que as réplicas ficam presas em ótimos locais a taxa de troca não é uma propriedade da escada (varia 0,10 entre sementes com a escada fixa, no TSP a partir de $T = 1$) e sobe ao longo da execução, à medida que a busca congela. Nesse caso a calibração acerta apenas na média; as opções são começar de uma primeira temperatura mais alta ou usar uma atualização dinâmica durante a execução.
\item \textbf{Tipos 3, 4 e 6 não são reprodutíveis com semente}, porque um par lê estatísticas de outros pares que podem estar em execução. O resultado é válido, mas duas execuções com a mesma semente podem diferir.
\item \textbf{Tipos 1 e 5 exigem escada crescente} (distribuições 1, 3 e 4). A combinação da distribuição 2 (decrescente) com as atualizações dinâmicas não foi tratada.
\item \textbf{Nenhuma regra dinâmica superou de forma consistente uma boa escada fixa} nos problemas testados. O valor prático delas está em não exigir essa boa escada de partida.
\item \textbf{A retomada de um checkpoint não é uma repetição bit a bit} (Seção 6.6).
\item \textbf{As probabilidades dos movimentos são fixas}; a API mede o uso de cada movimento, mas não adapta os pesos.
\item \textbf{O intervalo de atualização automático não melhora a qualidade} (Seção 9.5); elimina um parâmetro, e só.
\item \textbf{A construção em paralelo exige código de construção seguro para threads} e, havendo soluções iniciais de mesmo custo, a melhor solução da rodada zero pode ser qualquer uma delas.
\item \textbf{As execuções de um experimento são sequenciais}; o gerenciador não distribui execuções entre máquinas.
\item \textbf{O cenário do irace foi validado com orçamento pequeno} (300 execuções); os parâmetros que ele devolveu não são uma recomendação.
\item \textbf{As linhas de SD do benchmark não foram repetidas} depois das alterações posteriores à primeira medição. \end{itemize}

Ideias lidas e não implementadas: ajuste do número de réplicas durante a execução (Łącki e Miasojedow, 2016); PT em duas dimensões, com uma segunda escada sobre a penalidade das restrições (Delacour et al., 2025; Sim, Choi e Park, 2026); caminhos otimizados entre a distribuição de referência e a alvo (Syed et al., 2021); otimização das temperaturas por gradiente (Miyata, Arai e Takabe, 2026).

\section{Origem de cada método}\label{s12}

Resumo das atribuições feitas ao longo do documento. «Própria» indica uma regra ou adaptação criada nesta versão, sem publicação que a descreva.

\begin{tabela}{>{\raggedright\arraybackslash}p{0.25\linewidth}>{\raggedright\arraybackslash}p{0.31\linewidth}>{\raggedright\arraybackslash}X}
\toprule
\cab{Recurso} & \cab{Fonte} & \cab{O que é próprio da v2} \\
\midrule
PT e critério de troca & Geyer (1991); Hukushima e Nemoto (1996) & — \\
Modelo \textit{dataflow} (MC, SW, TU) & Tese do autor da API & Escalonamento (fila única, contadores atômicos, continuação na thread, espera em duas fases) \\
Tipo 1: 23\% & Atchadé, Roberts e Rosenthal (2011); equação da tese & Extremos fixos, correção nos dois sentidos, faixa morta, amortecimento \\
Tipo 2: taxas iguais & Hukushima (1999) & Proteção de ordem \\
Tipo 3: \textit{feedback} & Katzgraber et al. (2006), como na tese & Apenas a correção da condição de corrida \\
Tipo 4: \textit{feedback} por quantis & Katzgraber et al. (2006) (densidade ótima) & Cálculo global por quantis, extremos fixos, meio passo \\
Tipo 5: aproximação estocástica & Miasojedow, Moulines e Vihola (2013) & Expoente 0,6 e limite na temperatura máxima inicial \\
Tipo 6: rejeição igual & Syed et al. (2019) & Interpolação linear em 1/$T$, meio passo, ajuste durante a execução \\
Esquema par-ímpar & Okabe et al. (2001); Syed et al. (2019) & — \\
Idas e voltas & Katzgraber et al. (2006) & — \\
Barreira de comunicação & Syed et al. (2019) & — \\
\code{calibrateTemperatures} & Ideia relacionada: Ben-Ameur (2004) & Própria (amostra da segunda metade de uma descida; alvos padrão) \\
\code{calibrateLadder} & Alvo de 23\%: Atchadé, Roberts e Rosenthal (2011) & Regra de correção: Miasojedow, Moulines e Vihola (2013). Própria: pilotos do PT, média na segunda metade, limite de 90\% \\
Intervalo de atualização automático & Syed et al. (2019) (rodadas de ajuste dobradas) & Aplicação às seis regras da API \\
Integração com o irace & López-Ibáñez et al. (2016) & Opções de linha de comando e os três custos \\
Comparação de configurações & Mann e Whitney (1947); Wilcoxon (1945); Friedman (1937); Holm (1979); metodologia de Demšar (2006) & Valores-$p$ exatos com empates; tabelas \\
Tempo até o alvo & Aiex, Resende e Ribeiro (2007) & Apenas os dados do gráfico \\
Checkpoint, movimentos com probabilidades, critérios de parada, semente, estatísticas, tempo por componente, gerenciador de experimentos (curvas, partes, relatórios), construção por réplica & — & Recursos de engenharia desta versão \\
\bottomrule
\end{tabela}

\section*{Referências}\label{s13}\addcontentsline{toc}{section}{Referências}

\begin{referencias}
\refitem AIEX, R. M.; RESENDE, M. G. C.; RIBEIRO, C. C. TTT plots: a perl program to create time-to-target plots. \textit{Optimization Letters}, v. 1, p. 355–366, 2007.

\refitem ATCHADÉ, Y. F.; ROBERTS, G. O.; ROSENTHAL, J. S. Towards optimal scaling of Metropolis-coupled Markov chain Monte Carlo. \textit{Statistics and Computing}, v. 21, p. 555–568, 2011.

\refitem BEN-AMEUR, W. Computing the initial temperature of simulated annealing. \textit{Computational Optimization and Applications}, v. 29, n. 3, p. 369–385, 2004.

\refitem DELACOUR, C.; SAJEEB, M. M. H.; HESPANHA, J. P.; CAMSARI, K. Y. Two-dimensional parallel tempering for constrained optimization. arXiv:2506.14781, 2025. \textit{(lido; não implementado)}

\refitem DEMŠAR, J. Statistical comparisons of classifiers over multiple data sets. \textit{Journal of Machine Learning Research}, v. 7, p. 1–30, 2006.

\refitem FRIEDMAN, M. The use of ranks to avoid the assumption of normality implicit in the analysis of variance. \textit{Journal of the American Statistical Association}, v. 32, n. 200, p. 675–701, 1937.

\refitem GEYER, C. J. Markov chain Monte Carlo maximum likelihood. In: \textit{Computing Science and Statistics: Proceedings of the 23rd Symposium on the Interface}. Interface Foundation of North America, 1991.

\refitem HOLM, S. A simple sequentially rejective multiple test procedure. \textit{Scandinavian Journal of Statistics}, v. 6, n. 2, p. 65–70, 1979.

\refitem HUKUSHIMA, K. Domain-wall free energy of spin-glass models: numerical method and boundary conditions. \textit{Physical Review E}, v. 60, p. 3606–3613, 1999.

\refitem HUKUSHIMA, K.; NEMOTO, K. Exchange Monte Carlo method and application to spin glass simulations. \textit{Journal of the Physical Society of Japan}, v. 65, n. 6, p. 1604–1608, 1996.

\refitem LÓPEZ-IBÁÑEZ, M.; DUBOIS-LACOSTE, J.; PÉREZ CÁCERES, L.; BIRATTARI, M.; STÜTZLE, T. The irace package: iterated racing for automatic algorithm configuration. \textit{Operations Research Perspectives}, v. 3, p. 43–58, 2016.

\refitem KATZGRABER, H. G.; TREBST, S.; HUSE, D. A.; TROYER, M. Feedback-optimized parallel tempering Monte Carlo. \textit{Journal of Statistical Mechanics: Theory and Experiment}, v. 2006, n. 03, p. P03018, 2006.

\refitem ŁĄCKI, M. K.; MIASOJEDOW, B. State-dependent swap strategies and automatic reduction of number of temperatures in adaptive parallel tempering algorithm. \textit{Statistics and Computing}, v. 26, p. 951–964, 2016. arXiv:1406.3218. \textit{(lido; não implementado)}

\refitem MANN, H. B.; WHITNEY, D. R. On a test of whether one of two random variables is stochastically larger than the other. \textit{The Annals of Mathematical Statistics}, v. 18, n. 1, p. 50–60, 1947.

\refitem MIASOJEDOW, B.; MOULINES, E.; VIHOLA, M. An adaptive parallel tempering algorithm. \textit{Journal of Computational and Graphical Statistics}, v. 22, n. 3, p. 649–664, 2013. arXiv:1205.1076.

\refitem MIYATA, T.; ARAI, S.; TAKABE, S. Refined gradient-based temperature optimization for the replica-exchange Monte-Carlo method. arXiv:2601.13542, 2026. \textit{(lido; não implementado)}

\refitem OKABE, T.; KAWATA, M.; OKAMOTO, Y.; MIKAMI, M. Replica-exchange Monte Carlo method for the isobaric–isothermal ensemble. \textit{Chemical Physics Letters}, v. 335, p. 435–439, 2001.

\refitem ROBERTS, G. O.; GELMAN, A.; GILKS, W. R. Weak convergence and optimal scaling of random walk Metropolis algorithms. \textit{The Annals of Applied Probability}, v. 7, n. 1, p. 110–120, 1997.

\refitem SIM, K.; CHOI, S.; PARK, J. Solving integer linear programming with parallel tempering. arXiv:2605.29366, 2026. \textit{(lido; não implementado)}

\refitem SYED, S.; BOUCHARD-CÔTÉ, A.; DELIGIANNIDIS, G.; DOUCET, A. Non-reversible parallel tempering: a scalable highly parallel MCMC scheme. arXiv:1905.02939, 2019. Publicado em \textit{Journal of the Royal Statistical Society: Series B}, v. 84, n. 2, p. 321–350, 2022.

\refitem SYED, S.; ROMANIELLO, V.; CAMPBELL, T.; BOUCHARD-CÔTÉ, A. Parallel tempering on optimized paths. In: \textit{Proceedings of the 38th International Conference on Machine Learning (ICML)}, 2021. arXiv:2102.07720. \textit{(lido; não implementado)}

\refitem WILCOXON, F. Individual comparisons by ranking methods. \textit{Biometrics Bulletin}, v. 1, n. 6, p. 80–83, 1945.

\refitem Tese de doutorado do autor da PTAPI (Universidade Federal de Ouro Preto), 2024: Capítulo 5, implementação paralela para CPU, e os três métodos de ajuste de temperatura da v1.
 \end{referencias}

{\small Os dados de periódico de Ben-Ameur (2004), Okabe et al. (2001), Miasojedow et al. (2013), Łącki e Miasojedow (2016), Syed et al. (2021, 2022), López-Ibáñez et al. (2016), Aiex et al. (2007), Demšar (2006), Friedman (1937), Holm (1979), Mann e Whitney (1947) e Wilcoxon (1945) foram citados de memória e devem ser conferidos antes de uso em publicação; os demais vêm da lista de referências da tese ou da primeira página dos arquivos baixados.\par}

\end{document}
